% This is part of Mes notes de mathématique
% Copyright (c) 2008-2018
%   Laurent Claessens, Carlotta Donadello
% See the file fdl-1.3.txt for copying conditions.

%+++++++++++++++++++++++++++++++++++++++++++++++++++++++++++++++++++++++++++++++++++++++++++++++++++++++++++++++++++++++++++
\section{Éléments généraux de topologie}
%+++++++++++++++++++++++++++++++++++++++++++++++++++++++++++++++++++++++++++++++++++++++++++++++++++++++++++++++++++++++++++

%---------------------------------------------------------------------------------------------------------------------------
\subsection{Définitions et propriétés de base}
%---------------------------------------------------------------------------------------------------------------------------

\begin{definition}		\label{DefTopologieGene}
Soit \( X \), un ensemble et \( \mT \), une partie de l'ensemble de ses parties qui vérifie les propriétés suivantes.
\begin{enumerate}
\item
  Les ensembles \( \emptyset \) et \( X \) sont dans \( \mT \),
\item
  Une union quelconque\footnote{Par «quelconque» nous entendons vraiment quelconque : c'est à dire indicée par un ensemble qui peut autant être \( \eN\) que \( \eR\) qu'un ensemble encore considérablement plus grand.} d'éléments de \( \mT\) est dans \( \mT\).
\item
  Une intersection \emph{finie} d'éléments de \( \mT\) est dans \( \mT\).
\end{enumerate}
Un tel choix \( \mT \) de sous-ensembles de \( X \) est une  \defe{topologie}{topologie} sur \( X \), et les éléments de \( \mT \) sont appelés des \defe{ouverts}{ouvert}. On dit aussi que \( (X,\mT) \) (voire simplement \( X \) lorsqu'il n'y a pas d'ambiguïté) est un  \defe{\href{http://fr.wikipedia.org/wiki/Espace_topologique}{espace topologique}}{espace topologique}.
\end{definition}

\begin{definition}		\label{DefFermeVoisinage}
Si \(X \) est un espace topologique, un sous-ensemble \( F \) de \( X \) est dit \defe{fermé}{fermé} si son complémentaire, \( F^c \), est ouvert.

Si \(a \in X\), on dit que \(V \subset X\) est un \defe{voisinage}{voisinage} de \(a\) s'il existe un ouvert \(\mO \in \mT\) tel que \(a \in \mO\) et \(\mO \subset V\).
\end{definition}

\begin{lemma}   \label{LemQYUJwPC}
    Union et intersection de fermés.
    \begin{enumerate}
        \item
            Une intersection quelconque de fermés est fermée.
        \item       \label{ItemKJYVooMBmMbG}
            Une union finie de fermés est fermée.
    \end{enumerate}
\end{lemma}

\begin{proof}
    Soit \( \{ F_i \}_{i\in I} \) est un ensemble de fermés; nous avons
    \begin{equation}
        \left( \bigcap_{i\in I}F_i \right)^c=\bigcup_{i\in I}F_i^c.
    \end{equation}
    Le membre de droite est une union d'ouverts, c'est donc un ouvert; donc l'intersection qui apparaît dans le membre de gauche est le complémentaire d'un ouvert: c'est donc un fermé.

    De la même manière, le complémentaire d'une union finie de fermés est une intersection finie de complémentaires de fermés, et est donc ouvert\footnote{Un bon exercice est d'écrire ces unions et intersections, pour se convaincre que ça fonctionne.}.
\end{proof}

Dans un espace topologique, nous avons une caractérisation très importante des ouverts.
\begin{theorem}		\label{ThoPartieOUvpartouv}
    Une partie d'un espace topologique est ouverte si et seulement si elle contient un voisinage\footnote{Définition~\ref{DefFermeVoisinage}.} ouvert de chacun de ses éléments.
\end{theorem}

\begin{proof}
    Soit \( X\) un espace topologique et \( A\subset X\). Le sens direct est évident : $A$ lui-même est un ouvert autour de $x\in A$, qui est inclus dans $A$.

Pour le sens inverse, nous supposons que \( A\) contienne un ouvert autour de chacun de ses points. Pour chaque $x\in A$, choisissons $\mO_x\subset A$ un ouvert autour de $x$. Alors,
\begin{equation}	\label{EqAUniondesOx}
	A=\bigcup_{x\in A}\mO_x:
\end{equation}
en effet, d'une part, $A\subset\bigcup_{x\in A}\mO_x$ parce que chaque élément $x$ de $A$ est dans le $\mO_x$ corrrespondant, par construction; et d'autre part, $\bigcup_{x\in A}\mO_x\subset A$ parce que chacun des $\mO_x$ est inclus à $A$.

Ainsi, $A$ est égal à une union d'ouverts, cela prouve que $A$ est un ouvert.
\end{proof}
Le lemme \ref{LemMESSExh} est une version particulière de celui-ci, pour l'espace topologique \( \eR \). Une autre application typique est la proposition~\ref{PropMMKBjgY} et le théorème~\ref{ThoESCaraB}.

%--------------------------------------------------------------------------------------------------------------------------- 
\subsection{Base de topologie}
%---------------------------------------------------------------------------------------------------------------------------

\begin{definition}[Base de topologie]   \label{DefQELfbBEyiB}
    Une famille \( \mB\) d'ouverts de \( X\) est une \defe{base de la topologie}{base!de topologie} de \( X\) si pour tout \( x\in X\) et pour tout voisinage \( V\) de \( x\), il existe \( A\in \mB\) tel que \( x\in A\subset V\).
\end{definition}

\begin{proposition} \label{PropMMKBjgY}
    Si \( \mB\) est une base de la topologie de \( X\) alors tout ouvert de \( X\) est une union d'éléments de \( \mB\).
\end{proposition}

\begin{proof}
    Soit \( \mO\) un ouvert de \( X\); pour chaque \( x\in\mO\) nous considérons un ouvert \( U(x)\) tel que \( x\in U(x)\subset \mO\) (possible par le théorème~\ref{ThoPartieOUvpartouv}). Nous prenons alors \( B(x)\in\mB\) tel que
    \begin{equation}
        x\in B(x)\subset U(x)\subset \mO.
    \end{equation}
    Alors nous avons \( \mO=\bigcup_{x\in \mO}B(x)\).
\end{proof}
Notons toutefois que nous sommes loin d'avoir une union dénombrable en général.

%---------------------------------------------------------------------------------------------------------------------------
\subsection{Quelques exemples}
%---------------------------------------------------------------------------------------------------------------------------

\subsubsection{Une première vague}
%///////////////////////////

\begin{example}\label{DefTopologieGrossiere}
  Pour un ensemble \( X \) quelconque, on considère l'ensemble \( \mT = \{ \emptyset; X\} \). Avec cet ensemble, on confère à \(X \) une structure d'espace topologique - même si elle nous apprend peu de choses\dots{} La topologie ainsi posée sur \(X \) est appelée \defe{topologie grossière}{topologie!grossière}.
\end{example}

\begin{example}\label{DefTopologieDiscrete}
  Pour un ensemble \( X \) quelconque, on considère l'ensemble \( \mT \) constitué de toutes les parties de \( X \). Avec cet ensemble, on confère à nouveau une structure d'espace topologique à \(X \); toutes les parties sont des ouverts, et aussi des fermés. La topologie ainsi posée sur \(X \) est appelée \defe{topologie discrète}{topologie!discrète}.
\end{example}

\begin{example} [Toutes les topologies d'un ensemble à 3 éléments]
  On pose \( X = \{1, 2, 3\} \). Alors on peut munir \( X \) de 29 topologies différentes\footnote{Remercions Erwann Aubry d'en avoir fourni la liste exhaustive!  \url{https://math.unice.fr/~eaubry/Enseignement/L3/rappelstopo.pdf}}; saurez-vous les retrouver toutes?
\end{example}

\subsubsection{Topologie engendrée}
%//////////////////////////

\begin{propositionDef}\label{DefTopologieEngendree}
  Soit \( X \) un ensemble, et \( \mT_0 \) un sous-ensemble de parties de \( X \), contenant \( \emptyset \) et \( X \). On construit alors l'ensemble \( \mT \) par
  \begin{equation}
    \label{EqTopologieEngendree}
    \mT = \bigl\{\bigcup_{\alpha \in A} \bigcap_{i=1}^{n_a} \mO_{\alpha, i} \tq \mO_{\alpha,i} \in \mT_0 \forall \alpha, \forall i \bigr\}.
  \end{equation}
  Alors \(\mT \) est une topologie sur \(X\), qu'on appelle \defe{topologie engendrée}{topologie!engendrée par une famille} par \( \mT_0 \).
\end{propositionDef}

\begin{proof}
  Par définition et construction, la famille \( \mT \) contient le vide, l'espace entier, et est stable par union quelconque. Reste à établir qu'elle est stable par intersection finie. Réalisons-le pour deux éléments de \( \mT \), et une récurrence sur le nombre d'éléemnts permet de conclure.

  Soient donc deux éléments de \( \mT \):
  \begin{equation}
    O_1 = \bigcup_{\alpha \in A} \bigcap_{i=1}^{n_a} \mO_{\alpha, i}\quad\text{ et } O_2 = \bigcup_{\beta \in B} \bigcap_{j=1}^{n_\beta} \mO_{\beta, j}.
  \end{equation}
  Un élément \( x \) est dans l'intersection de ces deux ensembles s'il existe \( \alpha \in A \) et \( \beta \in B \) tels que pour tout \(i \in \{1, \cdots, n_\alpha\} \) et tout \(j \in \{1, \cdots, n_\beta\} \), \( x \in \mO_{\alpha, i}\) et \(x \in \mO_{\beta, j} \). Ce qui revient à dire qu'il existe \((\alpha, \betaø \in A \times B \) tel que \( x \) se trouve à l'intersection de tous les \(  \mO_{\alpha, i}\) et \(\mO_{\beta, j} \). En posant \( n_{\alpha, \beta} = n_\alpha + n_\beta\), puis, pour \((\alpha, \beta)\), et pour \( k  \in \{1, \cdots, n_{\alpha,\beta}\} \):
  \begin{equation}
    \mO_{(\alpha, \beta), k} =
    \begin{cases}
      \mO_{\alpha,k}& \text{si } k \leq n_\alpha,\\
      \mO_{\beta,k - n_\alpha}& \text{si } k > n_\alpha,\\
    \end{cases}
  \end{equation}
il s'ensuit que
\begin{equation}
   O_1 \cap O_2 = \bigcup_{\gamma \in A\times B} \bigcap_{k=1}^{n_\gamma} \mO_{\gamma, k}.
\end{equation}
\end{proof}

\begin{remark}
  La famille composée des intersections finies d'éléments de \( \mT_0 \) forme une base de la topologie engendrée.
\end{remark}

\subsubsection{Exemple fondamental: topologie de \texorpdfstring{$\eR$}{R}}
\begin{definition}[Topologie réelle engendrée par les intervalles ouverts] \label{DefTopologieREngendree}
  Pour \( X = \eR \), considérons \( \mT_0 \) l'ensemble constitué de \( \emptyset \), de \( \eR \), et de tous les intervalles ouverts (de la forme \( \mathopen] a, b \mathclose[ \), avec \( a, b \in \eR \)). La topologie ainsi engendrée est alors la topologie usuelle sur \( \eR \).
\end{definition}

\begin{remark} \label{RemIntervallesBaseTopoR}
  \begin{enumerate}
  \item
    Les intervalles ci-dessus sont une base de la topologie de \( \eR \). En effet, l'intersection des intervalles \( \mathopen] a_i, b_i \mathclose[ \) (pour \( i = 1,\ldots, n (n \in \eN) \) et les \( a_i \) et \( b_i \) réels, est soit vide, soit \( \mathopen] \max\{a_i\}, \min\{b_i\}[ \).
  \item
    Pour tout intervalle ouvert \( \mathopen]a,b\mathclose[ \), il existe \( x \in \eR \) et \( \epsilon > 0 \) tels que cet intervalle soit \( \mathopen]x - \epsilon,x+\epsilon\mathclose[ \). Il suffit de prendre \( x = \frac{a+b} 2 \) et \( \epsilon = \frac{b-a}2 \). Ainsi, les intervalles de la forme \( \mathopen] x - \epsilon, x = \epsilon \mathclose[ \) forment, pour \( x \in \eR \) et \( \epsilon > 0 \), une base de la topologie usuelle.
  \end{enumerate}

\end{remark}

\begin{example}[Voisinages de réels]
  Pour qu'un ensemble soit un voisinage de \( x \in \eR \), il faut et il suffit qu'il contienne un intervalle ouvert contenant \( x \). Ainsi, les ensembles suivants sont des voisinages de \( 3 \) dans \( \eR\) :
  \begin{itemize}
  \item
    \(\mathopen] 1 , 5 \mathclose[\);
  \item
    \(\mathopen[ 0 , 10 \mathclose]\);
  \item
    \(\eR\);
  \end{itemize}
  tandis que les ensembles suivants ne sont pas des voisinages de \(3\) dans \(\eR\) :
  \begin{itemize}
  \item
    \(\mathopen] 1 , 3 \mathclose[\);
  \item
    \(\mathopen] 1 , 3 \mathclose]\);
  \item
    \(\mathopen[ 0 , 5 [\setminus\{ 3 \}\).
  \end{itemize}
\end{example}

\begin{lemma}  \label{LemSupOuvPas}
Le supremum d'une partie ouverte de \( \eR\) n'est pas dans l'ensemble (et n'est donc pas un maximum).
\end{lemma}

\begin{proof}
Soit $\mO$, un ensemble ouvert et $s$, son supremum. Si $s$ était dans $\mO$, on aurait un voisinage de $s$ (et donc un intervalle $\mathopen]s - \epsilon, s + \epsilon \mathclose[$ contenant $s$) et contenu dans $\mO$. Mais le point $s+\epsilon/2$ est alors à la fois dans $\mathopen]s - \epsilon, s + \epsilon \mathclose[ \subset \mO$ et plus grand que $s$, ce qui contredit le fait que $s$ soit un supremum de $\mO$.
\end{proof}

\begin{remark}
  L'hypothèse de finitude de l'intersection est indispensable en termes de topologie. L'exemple suivant :
  \begin{equation}
    \bigcap_{n=1}^{\infty}\mathopen] -\frac{1}{ n } , \frac{1}{ n } \mathclose[ = \{ 0 \},
  \end{equation}
  est une intersection (infinie) d'intervalles ouverts (aux sens de la définition~\ref{DefIntervallesReels} et de l'exemple~\ref{DefTopologieREngendree}), égale à un singleton, qui est en fait un intervalle fermé.
\end{remark}


\subsubsection{Topologie produit}
%/////////////////////////////
\begin{definition}[Produit d'espaces topologiques, thème~\ref{THEMEooYRIWooDXZnhX}]      \label{DefIINHooAAjTdY}
    Soient \( X_1\),\ldots, \( X_n\) des espaces topologiques. Leur \defe{produit}{produit!espaces topologiques}\index{topologie!produit} est l'ensemble
    \begin{equation}
        X=\prod_{i=1}^nX_i
    \end{equation}
    muni de la topologie engendrée par les produits \(A_1\times \cdots\times A_n\), avec \( A_i\in X_i \) ouverts de chacun des ensembles.
\end{definition}

\begin{definition}\label{DefTopologieRn}
  Lorsque tous les espaces précédents sont \( \eR \) munis de la topologie usuelle\footnote{Définition \ref{DefTopologieREngendree}}, on engendre la topologie usuelle sur \( \eR^n\).
\end{definition}

\subsubsection{Topologie induite}
%//////////////////////////

\begin{example} [Topologie induite] \label{DefVLrgWDB}
  Soit un espace topologique \( (X, \mT) \), et soit \( Y \subset X \). Alors on peut munir \( Y \) de la topologie constituée des \( Y \cap \mO \), pour \( \mO \in \mT \): c'est ce qu'on appelle la \defe{topologie induite}{topologie!induite}.
\end{example}

\begin{lemma}[\cite{MonCerveau}]        \label{LemBWSUooCCGvax}
    Soit \( (X,\tau_X)\) un espace topologique et \( S\subset X\), un fermé de \( X\) sur lequel nous considérons la topologie induite \( \tau_S\). Si \( F\) est un fermé de \( (S,\tau_S)\) alors \( F\) est fermé de \( (X,\tau_X)\).
\end{lemma}

\begin{proof}
    Nous savons que le complémentaire de \( F\) dans \( S\) est un ouvert de \( (S,\tau_S)\) : il existe un ouvert \( \Omega\in \tau_X\) tel que \( S\setminus F=S\cap \Omega\). Si maintenant nous considérons le complémentaire de \( S\) dans \( X\) nous avons
    \begin{equation}
        F^c=(S\setminus F)\cup (X\setminus S)=(S\cap \Omega)\cup S^c=(S\cap \Omega)\cup(S^c\cap \Omega)\cup S^c=\Omega\cup S^c.
    \end{equation}
    Vu que \( \Omega\) et \( S^c\) sont des ouverts de \( X\), l'union est un ouvert. Donc \( F^c\in \tau_X\) et \( F\) est un fermé de \( X\).
\end{proof}


Si \( A\) est un ouvert de \( X\), on pourrait croire que la topologie induite n'a rien de spécial. Il est vrai que \( B\) sera ouvert dans \( A\) si et seulement s'il est ouvert dans \( X\), mais certaines choses surprenantes se produisent tout de même. Par exemple, prendre la topologie induite de \( \eR\) vers un fermé de \( \eR\) donne des ouverts un peu inhabituels\dots

\begin{example}  \label{ExKYZwYxn}
    Considérons, pour un \(\epsilon > 0 \), l'intervalle ouvert \( I = \mathopen]1 - \epsilon, 1 + \epsilon\mathclose[\); et considérons le sous-ensemble \(A = \mathopen[ 0 , 1 \mathclose]\). L'ouvert que \( I \) induit dans \( A \) est
    \begin{equation}
        \tilde I =\{ x\in\mathopen[ 0 , 1 \mathclose]\tq x \in I \}=\mathopen] 1-\epsilon , 1 \mathclose].
    \end{equation}
    Oui, cela est \emph{ouvert} dans \( \mathopen[ 0 , 1 \mathclose]\). C'est d'ailleurs un des ouverts de la topologie induite de \( \eR\) sur \( \mathopen[ 0 , 1 \mathclose]\).

    Donc pour la topologie de \( \mathopen[ 0 , 1 \mathclose]\), tous les intervalles ouverts (au sens de la topologie) sont soit de la forme \( \mathopen]x - \delta, x + \delta \mathclose[\), ou bien \( \mathopen]x - \delta,1 \mathclose]\), ou encore \( \mathopen[0, x + \delta \mathclose[\).
\end{example}

%---------------------------------------------------------------------------------------------------------------------------
\subsection{Comparaison de topologies}
%---------------------------------------------------------------------------------------------------------------------------

\begin{definition}
  Soit \( X \) un ensemble, muni de deux topologies \( \mT \) et \( \mU \). On dit que la topologie \( \mT \) est \defe{plus fine}{topologie!comparaison} que \( \mU \) si tout ouvert de \( \mU \) est un ouvert de \( \mT \).
\end{definition}

\begin{example}
  \begin{enumerate}
  \item
    Sur \( X = \{ 1, 2, 3 \} \), la topologie constituée de \( \emptyset, \{1\}, \{1, 2 \}\) et \( X \) est plus fine que celle constituée de \( \emptyset, \{1, 2 \}\) et \( X \).
  \item
    La topologie discrète\footnote{Exemple~\ref{DefTopologieDiscrete}.} est la plus fine de toutes les topologies.
  \item
    La topologie grossière\footnote{Exemple~{DefTopologieGrossiere}.} est la moins fine de toutes les topologies.
  \end{enumerate}
\end{example}

\begin{definition}
  Deux topologies \( \mT \) et \( \mU \) sur un espace topologique sont dites \defe{équivalentes}{topologie!topologies équivalentes} si  \( \mT \) est plus fine que \( \mU \), et  \( \mU \) est plus fine que \( \mT \).
\end{definition}

%---------------------------------------------------------------------------------------------------------------------------
\subsection{Adhérence, fermeture, intérieur, point d'accumulation et isolé}
%---------------------------------------------------------------------------------------------------------------------------

\subsubsection{Intérieur}
%///////////////////////

\begin{definition}      \label{DEFooSVWMooLpAVZRInt}
    Soient un espace topologique \( X\) et une partie \( A\) de \( X\).
    \begin{enumerate}
        \item
            Un point \( x\in X\) est \defe{intérieur}{point intérieur} à \( A\) s'il est contenu dans un ouvert inclus à \( A\). L'ensemble des points intérieurs de \( A\) est noté $\Int(A)$.\nomenclature[T]{$\Int(A)$}{intérieur de \( A\)}
        \item
            L'\defe{intérieur}{intérieur} de \( A\), notée \( \mathring A\), est l'union de tous les ouverts de \( X\) contenus dans \( A\).
    \end{enumerate}
\end{definition}
\begin{remark}\label{RemIntOuvert}
Quelques remarques en vrac.
\begin{enumerate}
\item Pour tout \( A \subset X\), l'ensemble \( \mathring A\) est un ouvert, comme union quelconque d'ouverts. C'est même le plus grand ouvert contenu dans \( A \), puisque, si \( O \) est un ouvert inclus à \( A \), alors \( O \subset \mathring A\).

\item Par ailleurs, on a \( \mathring A = \Int A \): en effet, \( x \in \Int A \) si et seulement s'il existe un ouvert contenant \( x \) et inclus à \( A \), si et seulement si \( x \) est dans l'union de tous les ouverts contenus dans \( A \), si et seulement si \( x \in \mathring A \).

\item On a  \( \mathring A \subset A \), et \( \mathring A = A \) si et seulement si $A$ est un ouvert: en sens direct, c'est clair par égalité d'ensembles; en sens inverse, c'est aussi clair puisque $A$ est alors un ouvert contenu à $A$, donc  \( A \subset \mathring A \).

\end{enumerate}
\end{remark}

\begin{example}
  Trouver l'intérieur d'un intervalle dans $\eR$ consiste à «ouvrir là où c'est fermé».
  \begin{enumerate}
  \item
    $\Int(\mathopen[ 0 , 1 \mathclose[)=\mathopen] 0 , 1 \mathclose[$.

    En effet, $\mathopen] 0,1  \mathclose[\subset\Int(\mathopen[ 0 , 1 \mathclose[)$, puisque c'est un ouvert. Réciproquement, montrons que $\Int( \mathopen[ 0 , 1\mathclose [ )\subset\mathopen] 0 , 1 \mathclose[$. Vu que l'intérieur d'un ensemble est inclus dans l'ensemble, nous savons déjà que $\Int( \mathopen[ 0 , 1\mathclose [ )\subset\mathopen[ 0 , 1 \mathclose[$. Nous devons donc seulement montrer que $0$ n'est pas dans l'intérieur de $\mathopen[ 0 , 1 \mathclose[$. C'est le cas parce que toute boule du type $B(0,r)$ contient le point $-r/2$ qui n'est pas dans $\mathopen[ 0 , 1 \mathclose[$.
  \item
    $\Int( \mathopen[ 0 , + \infty [)=\mathopen] 0 , + \infty \mathclose[$.
  \item
    $\Int( \mathopen] - \infty, 4 [ )=\mathopen] - \infty, 4 \mathclose[$.
  \item
    $\Int( \mathopen] 2 , 3 \mathclose[ )=\mathopen] 2 , 3 \mathclose[$.
  \end{enumerate}
\end{example}

\begin{proposition}\label{IntQestVide}
  On a \( \Int \eQ = \emptyset \).
\end{proposition}
\begin{proof}
  Si \( x, y \in \eQ \), alors le lemme~\ref{IntQestVide_corpsOrdonnes} dit que l'intervalle \( \mathopen] x, y \mathclose[ \) n'est pas dans \( \eQ \). Donc \( \eQ \) ne contient aucun ouvert de \( \eR \): son intérieur est donc vide. 
\end{proof}

\subsubsection{Adhérence et fermeture}
%///////////////////////

Disons-le tout de suite : «adhérence» et «fermeture» sont synonymes.
\begin{definition}      \label{DEFooSVWMooLpAVZR}
    Soient un espace topologique \( X\) et une partie \( A\) de \( X\).
    \begin{enumerate}
        \item
            Un point \( x\in X\) est \defe{adhérent}{point adhérent} à \( A\) si tout ouvert de \( X\) contenant \( x\) a une intersection non vide avec \( A\). L'ensemble des points d'adhérence de \( A\) est noté $\Adh(A)$.\nomenclature[T]{$\Adh(A)$}{adhérence de \( A\)}
        \item
            L'\defe{adhérence}{adhérence} de \( A\), notée \( \bar A\), est l'intersection de tous les fermés de \( X\) contenant \( A\).
    \end{enumerate}
\end{definition}

\begin{example}
	Dans $\eR$, l'infimum et le supremum d'un ensemble sont des points adhérents. En effet si $M$ est le supremum de $A\subset\eR$, pour tout $\varepsilon$, il existe un $a\in A$ tel que $a>M-\varepsilon$, tandis que $M>a$. Donc, pour tout $\varepsilon$, l'intervalle \( \mathopen ] M - \varepsilon, M + \varepsilon \mathclose[ \) a une intersection non-vide avec $A$.

	Le même raisonnement montre que l'infimum est également dans l'adhérence de $A$.
\end{example}

\begin{lemma}       \label{LEMooILNCooOFZaTe}
    L'adhérence de \( A\) est l'ensemble des points adhérents :
    \begin{equation}
        \Adh(A)=\bar A.
    \end{equation}
    Par ailleurs, on a le lien
    \begin{equation} \label{PropCompleIntBar}
      (\mathring A)^c = \widebar{A^c}.
    \end{equation}
\end{lemma}

\begin{proof}
    Commençons par prouver la dernière égalité d'ensembles. On a les équivalences entre les éléments suivants, pour tout $x \in X$:
    \begin{itemize}
    \item $x$ n'est pas dans $\mathring A$;
    \item il n'y a aucun ouvert contenant $x$ et inclus à $A$;
    \item tout ouvert contenant $x$ a une intersection non-vide avec $A^c$;
    \item $x$ est dans $\widebar{A^c}$.
    \end{itemize}
    Nous allons à présent montrer l'égalité d'ensembles \( \Adh(A)=\bar A \) en prouvant la double inclusion per contraposée.
    \begin{subproof}
        \item[Si \( x\in \bar A\) alors \( x\in\Adh(A)\)]
            Si \( x\) n'est pas dans \( \bar A\) alors nous avons un fermé \( F\) contenant \( A\) et pas \( x\). Le complémentaire \( F^c\) est un ouvert qui contient \( x\) et dont l'intersection avec \( A\) est vide. Donc \( x\) n'est pas dans \( \Adh(A)\).

        \item[Si \( x\in\bar A\) alors \( x\in \Adh(A)\)]

            Si \( x\) n'est pas dans \( \Adh(A)\) alors il existe un ouvert \( \mO\) contenant \( x\) et n'intersectant pas \( A\). Le complémentaire \( \mO^c\) est un fermé qui contient \( A\) et qui ne contient pas \( x\).

            Vu que \( \bar A\) est l'intersection de tous les fermés contenant \( A\), nous avons \( \bar A\subset\mO^c\) et donc \( x\) n'est pas dans \( \bar A\).
    \end{subproof}
\end{proof}

\begin{remark}\label{RemAdhFerme}\label{EXooICLBooJzQFNY}
  Comme corollaire du lemme précédent, et grâce aux remarques faites pour les intérieurs, on obtient que pour \( A \subset X \) :
  \begin{enumerate}
  \item
    l'ensemble \( \bar A \) est fermé: c'est en effet le complémentaire d'un ouvert, précisément l'intérieur de \( A^c \);
  \item
    l'ensemble $\bar A$ est le plus petit fermé contenant $A$. Cela signifie que si $B$ est un fermé qui contient $A$, alors $\bar A\subset B$. Cela est fondamentalement le sens de la définition~\ref{DEFooSVWMooLpAVZR}.
  \item
    clairement, \( A \subset \bar A \); de plus, \( A \) est fermé si et seulement si \( \bar A = A \): en effet, \( A \) est fermé si et seulement si \( A^c \) est ouvert, si et seulement si l'intérieur de \( A^c \) est \( A^c \) lui-même; or, l'intérieur de \( A^c \) est le complémentaire de \( \bar A \) par le lemme \ref{LEMooILNCooOFZaTe}, si bien que \( A \) est fermé si et seulement si \( (\bar A)^c  = A^c \), ou encore\dots{} ce qu'on affirmait au début.
  \end{enumerate}
\end{remark}

\begin{example}
  Trouver l'adhérence d'un intervalle dans $\eR$ consiste à «fermer là où c'est ouvert», en prenant garde de ne pas fermer en les infinis.
 \begin{enumerate}
  \item
    $\Adh\big(\mathopen[ 0 , 1 [\big)=\mathopen[ 0 , 1 \mathclose]$.

    En effet, d'abord $\Adh(\mathopen[ 0 , 1 \mathclose[ )\subset\mathopen[ 0,1  \mathclose])$, puisque $\mathopen[ 0,1  \mathclose]$ est un fermé: c'est le complémentaire dans $\eR$ de $\mathopen]-\infty, 0\mathclose [ \cup \mathopen]1, +\infty\mathclose[$. Réciproquement, montrons que $\mathopen[ 0,1  \mathclose] \subset \Adh(\mathopen[ 0 , 1 \mathclose[)$.  Vu que tout ensemble est inclus dans son adhérence, nous savons déjà que $\mathopen[ 0 , 1 [\subset\Adh( \mathopen[ 0 , 1 [ )$. Nous devons donc seulement montrer que $1$ est aussi dans l'adhérence de $\mathopen[ 0 , 1 [$. C'est le cas parce que toute boule du type $B(1,r)$ contient le point $1-r/2$ qui est dans $\mathopen[ 0 , 1 [$.
  \item
    $\Adh( \mathopen[ 0 , + \infty [)=\mathopen[ 0 , + \infty \mathclose[$. Remarquez que le crochet à $+\infty$ reste ouvert, et pour cause: $+\infty \notin \eR$.
  \item
    $\Adh( \mathopen] - \infty, 4 [ )=\mathopen] - \infty, 4 \mathclose]$.
  \item
    $\Adh( \mathopen[ 2 , 3 \mathclose] )=\mathopen[ 2 , 3 \mathclose]$.
  \end{enumerate}
\end{example}

\begin{remark}
  La définition \ref{DefIntervallesReels} prend alors tout son sens quant aux intervalles ouverts et fermés: ils sont à juste titre respectivement ouverts et fermés pour la topologie usuelle.
\end{remark}

\begin{remark}
  Attention à ne pas confondre $(\bar A)^c$ et
  $\overline{ A^c }$. Ces deux ensembles ne sont, en général, pas
  égaux. En effet, en tant que complément d'un fermé, l'ensemble
  $(\bar A)^c$ est certainement ouvert, tandis que, en tant
  que fermeture, l'ensemble $\overline{ A^c }$ est
  fermé. Pouvez-vous trouver des exemples d'ensembles $A$ tels que
  $(\bar A)^c=\overline{ A^c }$ ?
\end{remark}


\begin{example}		\label{ParlerEncoredeF}
	Il ne faut pas conclure de l'exemple précédent qu'un point adhérent est automatiquement un minimum ou un maximum. En effet, si nous regardons l'ensemble formé par les points de la suite $x_n=(-1)^n/n$, le nombre zéro est un point adhérent (et une limite, nous définirons cela plus tard) mais il n'est pas un infimum ni un maximum.
\end{example}

\begin{proposition}
	Soient $A$ et $B$ deux parties de l'espace topologique $X$.
	\begin{description}
		\item [Inclusions:]
                  si $A\subset B$, alors $\mathring A \subset \mathring B$ et $\bar A\subset\bar B$.
		\item [Unions:]
                  $\overline{ A\cup B }=\overline{ A }\cup\overline{ B }$ et $\mathring A \cup \mathring B \subset \Int(A\cup B)$.
		\item [Intersections:] $\mathring A \cap\mathring B =\Int(A\cap B)$ et $\overline{ A\cap B }\subset\bar A\cap\bar B$.
	\end{description}
\end{proposition}

\begin{proof}
	\begin{enumerate}
		\item
                  Si $a$ est dans l'intérieur de $A$, il existe un ouvert autour de $a$ contenu dans $A$. Cet ouvert est alors contenu dans $B$ et donc $a$ est dans l'intérieur de $B$. Si maintenant $a$ est dans l'adhérence de $A$, tout ouvert contenant $a$ contient un élément de $A$ et donc un élément de $B$, ce qui prouve que $a$ est dans l'adhérence de $B$.
		\item
                  Nous avons $A\subset A\cup B$ et donc, en utilisant le premier point, $\bar A\subset\overline{ A\cup B }$. De la même manière, $\bar B\subset\overline{ A\cup B }$. En prenant l'union, $\bar A\cup\bar B\subset\overline{ A\cup B }$. De la même façon, on obtient $\mathring A \cup \mathring B \subset \Int(A\cup B)$.

                  Montrons l'inclusion $\overline{ A\cup B }\subset \bar A\cup\bar B$ par contraposée. Supposons que $a$ ne soit ni dans $\bar A$ ni dans $\bar B$. Il existe donc des ouverts $\mO_A$ et $\mO_B$ contenant tous les deux $a$, et tels que
			\begin{equation}
				\begin{aligned}[]
					\mO_A\cap A&=\emptyset,\\
					\mO_B \cap B&=\emptyset.
				\end{aligned}
			\end{equation}
                        On considère alors $\mO = \mO_A \cap \mO_B$: c'est un ouvert non-vide, puisque $a$ est dedans; et il n'intersecte ni $A$, ni $B$, donc $\mO \cap (A \cup B) = \emptyset$, et $a\notin\overline{ A\cup B }$.

		\item
			Si nous appliquons le second point à $A^c$ et $B^c$, nous trouvons
			\begin{equation}
				\overline{ A^c\cup B^c }=\overline{A^c}\cup\overline{ B^c}.
			\end{equation}
			En utilisant les propriétés du lemme~\ref{LemPropsComplement}, le membre de gauche devient
			\begin{equation}	\label{Eq2707CACBCAB}
				\overline{ A^c\cup B^c }=\overline{(A\cap B)^c }=(\Int(A\cap B))^c,
			\end{equation}
			tandis que le membre de droite devient
			\begin{equation}		\label{Eq2707cAcBACAACB}
				\overline{ A^c }\cup\overline{ B^c }=(\mathring A)^c\cup(\mathring B)^c=\Big( \Int(A)\cap\Int(B) \Big)^c.
			\end{equation}
			En égalant le membre de droite de \eqref{Eq2707CACBCAB} avec celui de \eqref{Eq2707cAcBACAACB} et en passant au complémentaire nous trouvons
			\begin{equation}
				\Int(A\cap B)=\mathring A \cap\mathring B,
			\end{equation}
			comme annoncé.

			La dernière affirmation provient du fait que $\mathring A \subset \Int(A\cup B)$ et de la propriété équivalente pour $B$.
	\end{enumerate}
\end{proof}

\begin{remark} \label{ExBFLooUNyvbw}
	Nous avons prouvé que $\overline{ A\cap B }\subset\bar A\cap\bar B$. Il arrive que l'inclusion soit stricte, comme dans l'exemple suivant. Si nous prenons $A=\mathopen[ 0 , 1 \mathclose]$ et $B=\mathopen] 1 , 2 \mathclose]$, nous avons $A\cap B=\emptyset$ et donc $\overline{ A\cap B }=\emptyset$. Par contre nous avons $\bar A\cap\bar B=\{ 1 \}$.

        Le lecteur est invité à montrer que \( \mathring A \cup \mathring B \neq \Int(A \cup B) \) pour les mêmes ensembles.
\end{remark}


\begin{lemma}       \label{LemkUYkQt}
    Si \( B\subset A\) alors la fermeture de \( B\) pour la topologie de \( A\) (induite de \( X\)) que nous noterons \( \tilde B\) est
    \begin{equation}
        \tilde B=\bar B\cap A
    \end{equation}
    où \( \bar B\) est la fermeture de \( B\) pour la topologie de \( X\).
\end{lemma}

\begin{proof}
    Si \( a\in \bar B\cap A\), un ouvert de \( A\) autour de \( a\) est un ensemble de la forme \( \mO\cap A\) où \( \mO\) est un ouvert de \( X\). Vu que \( a\in\bar B\), l'ensemble \( \mO\) intersecte \( B\) et donc \( (\mO\cap A)\cap B\neq \emptyset\). Donc \( a\) est bien dans l'adhérence de \( B\) au sens de la topologie de \( A\).

    Pour l'inclusion inverse, soit \( a\in \tilde  B\), et montrons que \( a\in \bar B\cap A\). Par définition \( a\in A\), parce que \( \tilde B\) est une fermeture dans l'espace topologique \( A\). Il faut donc seulement montrer que \( a\in\bar B\). Soit donc \( \mO\) un ouvert de \( X\) contenant \( a\); par hypothèse \( \mO\cap A\) intersecte \( B\) (parce que tout ouvert de \( A\) contenant \( a\) intersecte \( B\)). Donc \( \mO\) intersecte \( B\). Cela signifie que tout ouvert (de \( X\)) contenant \( a\) intersecte \( B\), ou encore que \( a\in \bar B\).
\end{proof}

Adhérence et topologie induite peuvent créer quelques surprises, ainsi que le montre l'exemple suivant.

\begin{example} \label{ExloeyoR}
Prenons \( X=\eR\) et \( A=\mathopen] 0 , 1 \mathclose[\). Si \( B=\mathopen] \frac{ 1 }{2} , 1 \mathclose[ \), alors la fermeture de \( B\) dans \( A\) sera \( \tilde B=\mathopen[ \frac{ 1 }{2} , 1 [\) et non \( \mathopen[ \frac{ 1 }{2} , 1 \mathclose]\) comme on l'aurait dans \( \eR\).
\end{example}

\begin{definition}\label{DefEnsembleDense}
  Soit \( X \) un espace topologique. Un sous-ensemble \( A \) de \( X \) est \defe{dense}{dense} dans \( X \) si \( \bar A = X\). 
\end{definition}

\begin{proposition}     \label{PropooUHNZooOUYIkn}
    Les rationnels sont denses dans les réels.
\end{proposition}
\index{densité!de \( \eQ\) dans \( \eR\)}

\begin{proof}
  Le lemme~\ref{LemooHLHTooTyCZYL} dit que, pour tous réels \( x, y \), il existe un rationnel dans \( \mathopen] x, y \mathclose[\). Donc tout intervalle ouvert contient un rationnel, et donc, \( \bar \eQ = \eR \).
\end{proof}

\subsubsection{Frontière}
%/////////////////////////
Comme pour tout ensemble \( A \) inclus à un espace topologique, on a \( \mathring A \subset A \subset \bar A \), on peut regarder les points entre l'adhérence et l'intérieur.
\begin{definition}	\label{PropDescFrpbsmI}
  Soit \( X \) un espace topologique, et \( A \subset X \). La \defe{frontière}{frontière} de \( A \), notée \( \partial A \), est l'ensemble des points adhérents de \( A \) qui ne sont pas intérieurs à \( A \). Ainsi,
  \begin{equation}
    \partial A = \bar A \setminus \mathring A.
  \end{equation}
\nomenclature[T]{$\partial A$}{La frontière de l'ensemble $A$}
\end{definition}

\begin{lemma}[Caractérisation équivalente de la frontière]      \label{LEMooEUYEooYcUfKr}
    Soient \( X\) un espace topologique et \( S\subset X\). Un point \( x\in X\) est dans \( \partial S\) si et seulement si tout voisinage de \( x\) contient un point de \( S\) et un point de \( S^c\).
\end{lemma}

\begin{proof}
    Supposons que tout voisinage de \( x\) contienne un point de \( S\) et un point de \( S^c\). Alors \( x\in \Adh(S)\) (définition~\ref{DEFooSVWMooLpAVZR}), mais pas dans l'intérieur de \( S\) parce que \( x\) ne possède pas de voisinage contenu dans \( S\). Donc \( x\in \partial S\).

    À l'inverse, si \( x\in\partial S\) alors \( x\) est dans l'adhérence de \( S\) et tout voisinage de \( x\) contient un point de \( S\). Mais \( x\) n'est pas dans l'intérieur de \( S\) et tout voisinage de \( x\) contient un point qui n'est pas dans \( S\), c'est-à-dire un point de \( S^c\).
\end{proof}

\begin{corollary}
    Un ensemble et son complémentaire ont même frontière.
\end{corollary}

\begin{proof}
    Conséquence du lemme~\ref{LEMooEUYEooYcUfKr}. Les points de \( \partial(S^c)\) sont caractérisés par le fait que tout voisinage d'un tel point contient un point de \( S^c\) et un point de \( (S^c)^c=S\).
\end{proof}

\begin{lemma}
  Soit \( A \subset X \), où \( X \) est un espace topologique. La frontière de $A$ peut également s'exprimer des façons suivantes :
	\begin{equation}
		\partial A= \bar A\cap(\mathring A)^c = \bar A\cap\overline{ A^c}.
	\end{equation}
\end{lemma}

\begin{proof}
	En partant de $\partial A=\bar A\setminus \mathring A$, la première égalité est une application de la propriété~\ref{ItemLemPropComplementiii} du lemme~\ref{LemPropsComplement}. La seconde égalité est alors~\eqref{PropCompleIntBar}.
\end{proof}

\begin{lemma}
  Soit \( A \subset X \), où \( X \) est un espace topologique. Alors \( \partial A = \emptyset \) si et seulement si \( A \) est à la fois ouvert et fermé.
\end{lemma}

\begin{proof}
  Si \( \partial A = \emptyset \), alors \( \bar A \subset \mathring A \); or on a aussi \( \mathring A \subset A \subset \bar A \). Les ensembles \( \mathring A \), \(A \) et \( \bar A \) sont donc tous égaux, ce qui, grâce aux remarques~\ref{RemIntOuvert} et \ref{RemAdhFerme}, est équivalent au fait que \( A \) soit à la fois ouvert et fermé.

  Réciproquement, si \( A \) est ouvert et fermé, alors \( \mathring A = \bar A \), et donc  \( \partial A = \emptyset \).
\end{proof}

\begin{example}
	Dans $\eR$, la frontière d'un intervalle borné est la paire constituée des points extrêmes. En effet
	\begin{equation}
		\partial\mathopen[ a , b [=\overline{ \mathopen[ a , b [ }\setminus\Int\big( \mathopen[ a , b [ \big)=\mathopen[ a , b \mathclose]\setminus\mathopen] a , b \mathclose[=\{ a,b \}.
	\end{equation}

	Toujours dans $\eR$ nous avons
	\begin{equation}
		\partial\eR=\bar\eR\setminus\Int(\eR)=\eR\setminus\eR=\emptyset,
	\end{equation}
	et, grâce aux propositions~\ref{IntQestVide} et \ref{PropooUHNZooOUYIkn},
	\begin{equation}
		\partial\eQ=\bar\eQ\setminus\Int(\eQ)=\eR\setminus\emptyset=\eR.
	\end{equation}
\end{example}

\begin{remark}
    Il serait toutefois faux de croire que \( \partial A=\partial \bar A\) pour toute partie \( A \subset X\). En effet, si \( A=\eR\setminus\{ 0 \}\) nous avons \( \partial A=\{ 0 \}\) et \( \bar A=\eR\), donc \( \partial \bar A=\emptyset\).
\end{remark}


\begin{remark}
    Petite remarque à propos des frontières. Soit \( X=\mathopen[ 0 , 1 \mathclose]\) muni de la topologie induite de \( \eR \). Les points \( 0\) et \( 1\) \emph{ne sont pas} dans la frontière de $X$. En effet, un ouvert autour de \( 1\) est un ensemble de la forme
    \begin{equation}
        \mathopen]1-r, 1+s \mathclose[ \cap X =\{ x\in X\tq 1 - r < x < 1 + s \}=\mathopen] 1-r , 1 \mathclose]
    \end{equation}
    où nous avons supposé \( 0<r<1\) et \( s > 0 \).

    Les points \( 0\) et \( 1\) sont par contre sur la frontière de \( \mathopen[ 0 , 1 \mathclose]\) lorsque cet ensemble est vu comme partie de l'espace topologique \( \eR\).
\end{remark}

\subsubsection{Points d'accumulation et isolés}
%/////////////////////////

\begin{definition}      \label{DEFooGHUUooZKTJRi}
    Soient un espace topologique \( X\) et une partie \( A\) de \( X\). Un point \( s\in X \) est un \defe{point d'accumulation}{point d'accumulation} de \( A\) si tout ouvert contenant \( s\) contient au moins un élément de \( A\setminus\{ s \}\).
\end{definition}

\begin{proposition}
  Tout point d'accumulation est un point d'adhérence.
\end{proposition}

\begin{proof}
  Il suffit de se référer aux définitions \ref{DEFooGHUUooZKTJRi} et \ref{DEFooSVWMooLpAVZR}. Si \( s \in X \) est d'accumulation pour un ensemble \( A \), alors tout ouvert le contenant rencontre \( A \setminus \{s \}\) et donc, rencontre \( A \) aussi.
\end{proof}

Quelle est la différence entre un point d'accumulation et un point d'adhérence ? La différence est que tous les points de \( A\) sont des points d'adhérence de \( A\), parce que tout voisinage de \( a\in A\) contient au moins \( a\) lui-même, alors que certains points de \( A\) peuvent ne pas être des points d'accumulation de \( A\).

\begin{example}     \label{EXooWOYQooJolaTV}
    Soit \( D=\mathopen] 1 , 2 \mathclose[\cup\{ 12 \}\). Le point \( 12\) est d'adhérence, mais pas d'accumulation parce que le voisinage \( \mathopen] 9 , 14 \mathclose[\) n'intersecte pas \( D\setminus \{ 12 \}\).

\end{example}

\begin{definition}      \label{DEFooXIOWooWUKJhN}
    Soient un espace topologique \( X\) et une partie \( A\) de \( X\). Un point \( s\in A \) est un \defe{point isolé}{point isolé} de \( A\) si il existe un voisinage ouvert \( \mO\) de \( s\) dans \( X\) tel que \( A\cap\mO=\{ s \}\).
\end{definition}

\begin{example}
	Soit $D=\{ \frac{1}{ n }\}_{n\in\eN}$. Tous les points de cet ensemble sont des points isolés (vérifier !).  Aucun point de $D$ n'est point d'accumulation. Cependant $0$ est un point d'accumulation.
\end{example}

\begin{example}
	Prenons $D=\mathopen[ 0 , 1 [\cup\mathopen] 2 , 3 \mathclose]$. Cet ensemble n'a pas de points isolés, et l'ensemble de ses points d'accumulation est $\mathopen[ 0 , 1 \mathclose]\cup\mathopen[ 2,3  \mathclose]$.

	Notez que les points $1$ et $2$ sont des points d'accumulation de $D$ qui ne font pas partie de $D$. Il est possible d'être un point d'accumulation de $D$ sans être dans $D$, mais pour être un point isolé dans $D$, il faut être dans $D$.
\end{example}

\begin{remark}
    À propos de la position des points d'accumulation et des points isolés.
    \begin{enumerate}
        \item
            Les points intérieurs sont tous des points d'accumulation.
        \item
            Les points isolés ne sont jamais intérieurs.\quext{Ça, seulement si la topologie est bien foutue. Imaginez qu'un point soit ouvert à lui tout seul\dots}
        \item
            Certains points d'accumulation ne font pas partie de l'ensemble. Par exemple le point $1$ est un point d'accumulation de $E=\mathopen] 0 , 1 \mathclose[$.
        \item
            Les points de la frontière sont soit d'accumulation soit isolés.
    \end{enumerate}
\end{remark}

\begin{example}
	Tous les points de $\eR$ sont des points d'accumulation de $\eQ$ parce que dans toute boule autour d'un réel, on peut trouver un nombre rationnel.
\end{example}

\begin{remark}
	L'ensemble des points d'accumulation d'un ensemble n'est pas exactement son adhérence. En effet, un point isolé dans $A$ est dans l'adhérence de $A$, mais n'est pas un point d'accumulation de $A$.
\end{remark}

Nous voyons maintenant ce que ces définitions donnent dans le cas de l'espace topologique \( \eR\).

\begin{lemma}
    Un point $a\in \eR$ est un point d'accumulation de $D$ si pour tout $\varepsilon>0$,
    \begin{equation}
        \Big( \mathopen] a-\varepsilon , a+\varepsilon \mathclose[\setminus\{ a \} \Big)\cap D\neq\emptyset.
    \end{equation}
    Autrement dit, quel que soit l'intervalle autour de  $a$ que l'on considère, le point $a$ n'est pas tout seul dans $D$.
\end{lemma}

\begin{lemma}
    Soit $D\subset\eR$. Un point $a\in D$ est isolé dans $D$ si et seulement si existe $\varepsilon>0$ tel que
    \begin{equation}
        \mathopen] a-\varepsilon , a+\varepsilon \mathclose[\cap D=\{ a \}.
    \end{equation}
    Autrement dit, il existe un intervalle autour de $a$ dans lequel $a$ est le seul élément de $D$.
\end{lemma}

%---------------------------------------------------------------------------------------------------------------------------
\subsection{Convergence de suites}
%---------------------------------------------------------------------------------------------------------------------------

Dès que nous avons une topologie nous avons une notion de convergence.
\begin{definition}[Convergence de suite] \label{DefXSnbhZX}
    Une suite $(x_n)$ d'éléments de $E$ \defe{converge}{convergence!de suite} vers un élément $y$ de $E$ si pour tout ouvert $\mO$ contenant $y$, il existe un $K\in \eN$ tel que $k>K$ implique $x_k\in\mO$.
\end{definition}
\index{limite!de suite!espace topologique}


\begin{proposition}[\cite{MonCerveau}]      \label{PROPooBBNSooCjrtRb}
    Une suite contenue dans un fermé ne peut converger que vers un élément de ce fermé.
\end{proposition}

\begin{proof}
    Soient un espace topologique \( X\) et un fermé \( F\) dans \( X\). Nous supposons que la suite \( (x_k)\) soit contenue dans \( F\). Nous allons prouver qu'aucun élément de \( F^c\) ne peut être limite.

    Soit \( a \in F^c\). Vu que le complémentaire de \( F\) est un ouvert, et vu le théorème \ref{ThoPartieOUvpartouv}, il existe un ouvert \( \mO_a\) contenant \( a\), et contenu dans \( F^c\). Le voisinage \( \mO_a\) de \( a\) ne contient donc aucun élément de la suite \( (x_k)\), qui ne peut donc pas converger vers \( a\).
\end{proof}

\begin{corollary}\label{CorLimAbarA}
  Soit \( A \) un sous-ensemble d'un espace topologique \(X \). Toute suite d'éléments de \(A \) qui converge, admet pour limite un élément de \( \bar A \).
\end{corollary}
\begin{proof}
  Une fois la suite \( (x_n) \) fixée, il suffit de remarquer que tous les \( x_n \) sont dans \( \bar A \), et puis d'appliquer la proposition~\ref{PROPooBBNSooCjrtRb}. 
\end{proof}


\begin{lemma}   \label{LemPESaiVw}
    Soit \( A\subset X\) muni de la topologie induite de \( X\) et \( (x_n)\) une suite dans \( A\). Si \( (x_n) \) converge vers un élément \( x \) dans \(A \), alors elle converge aussi vers \(x \) dans \( X \).
\end{lemma}

\begin{proof}
    Soit \( \mO\) un ouvert autour de \( x\) dans \( X\). Alors \( A\cap\mO\) est un ouvert autour de \( x\) dans \( A\) et il existe \( N\in \eN\) tel que si \( n\geq N\), alors \( x_n\in A\cap\mO\subset\mO\).
\end{proof}

%---------------------------------------------------------------------------------------------------------------------------
\subsection{Pour des limites uniques : séparabilité}
%---------------------------------------------------------------------------------------------------------------------------

Notons que l'on a parlé d'\emph{une} limite de suite jusqu'à présent: en effet, s'il existe deux éléments distincts $x$ et $y$ tels que tout ouvert contenant $x$ contient $y$, alors la définition \ref{DefXSnbhZX} dit que toute suite convergeant vers $y$ converge aussi vers $x$\dots


\begin{example} \label{EXooSHKAooZQEVLB}
    Oui, il y a moyen de converger vers plusieurs points distincts si l'espace n'est pas super cool. Nous pouvons par exemple\cite{EJVQuas} considérer la droite réelle munie de sa topologie usuelle\footnote{Définition~\ref{DefTopologieREngendree}} et y ajouter un point $0'$ (qui clone le réel $0$) dont les voisinages sont les voisinages de $0$ dans lesquels nous remplaçons $0$ par $0'$. Dans cet espace, la suite $(1/n)$ converge à la fois vers $0$ et $0'$.

    En fait, on «voit» le problème: on ne peut pas distinguer d'un point de vue topologique le $0$ et le $0'$.
\end{example}

Nous posons la définition suivante, qui nous permettra de donner une assez grande classe d'espaces topologiques dans lesquels nous avons unicité de la limite\footnote{Voir la proposition \ref{PropFObayrf}.}.
\begin{definition}[Espace topologique séparé]  \label{DefYFmfjjm}\label{DefWEOTrVl}
    Si deux points distincts admettent toujours deux voisinages disjoints\footnote{Définition~\ref{DefEnsemblesDisjoints}.}, nous disons que l'espace est \defe{séparé}{espace!séparé} ou \defe{Hausdorff}{Hausdorff}.
\end{definition}

Attention, cette notion est à ne pas confondre avec :
\begin{definition}[Espace topologique séparable]  \label{DefUADooqilFK}
    Un espace topologique est \defe{séparable}{séparable!espace topologique} s'il possède une partie dénombrable\footnote{Définition~\ref{DefEnsembleDenombrable}.} dense\footnote{Définition~\ref{DefEnsembleDense}.}.
\end{definition}

\begin{proposition}\label{PropUniciteLimitePourSuites}
  Dans un espace séparé, si une suite converge, alors sa limite est unique.
\end{proposition}
\begin{proof}
  Supposons que la suite \( (x_k)\) converge vers deux éléments distincts \( x \) et \( y \). L'espace étant séparé, il existe deux ouverts \( \mO_x \) et \( \mO_y \), disjoints, contenant respectivement \( x \) et \( y \). La suite convergeant à la fois vers \( x \) et \( y \), il existe \( k_x \) et \( k_y \), tels que, si \( k \geq \max\{k_x, k_y\} \), l'élément  \( x_k \) est (à la fois) dans  \( \mO_x \) et \( \mO_y \). Cela est en contradiction avec le fait que ces deux ensembles sont disjoints.
\end{proof}

\begin{normaltext}
  Donc, on pourra parler, avec des espaces séparés, de «la limite d'une suite». On notera \( x_n\to a\), ou \(\lim_{n\to \infty} x_n = a \), pour signifier que la suite \( (x_n) \) converge vers \( a \). 
\end{normaltext}

\begin{proposition}[\cite{MonCerveau}]      \label{PROPooNRRIooCPesgO}
    La convergence d'une suite pour la topologie de l'espace produit implique la convergence des suites «composante par composante».
\end{proposition}

\begin{proof}
    Pour simplifier les notations, nous allons considérer le produit de deux espaces. Soit donc \( (x_k,y_k)\stackrel{X\times Y}{\longrightarrow}(x,y)\) et des ouverts \( \mO_1\) dans \( X\) autour de \( x\) et \( \mO_2\) autour de \( y\) dans \( Y\). Comme la partie \( \mO_1\times \mO_2\) est ouverte dans \( X\times Y\), il existe \( K\) tel que \( k>K\) implique \( (x_k,y_k)\in \mO_1\times \mO_2\). En particulier, quand \( k>K\), on a \( x_k\in \mO_1\) et  \( y_k\in \mO_2\). Par conséquent, \( x_k\stackrel{X}{\longrightarrow}x\) et \( y_k\stackrel{X}{\longrightarrow}y\). 
\end{proof}

\begin{proposition}[\cite{MonCerveau}]      \label{PROPooNRRIooCPesgO2}
  Soit \( X \) et \( Y \) des espaces topologiques. Soit \( (x_n, y_n) \) une suite de l'espace \( X \times Y \), muni de la topologie produit. Si \( x_n \to x \) dans \( X \) et \( y_n \to y \) dans \( Y \), alors la suite  \( (x_n, y_n) \) converge vers \( (x,y) \) dans \( X \times Y \).

  Cela s'étend à tout produit fini d'espaces topologiques.
\end{proposition}

\begin{proof}
  Soit \( V \) un voisinage de \( (x,y) \). Puisque \( X \times Y \) est muni de la topologie produit, il existe \( \mO_1 \) un ouvert de \( X \) et  \( \mO_2 \) un ouvert de \( Y \) tels que \(\mO_1 \times \mO_2 \subset V\). Vu les convergences des suites \( (x_n) \) et \( (y_n) \), il existe un entier naturel \( K_1 \) et un entier naturel \( K_2 \) tels que \( x_n \in \mO_1 \) quand \( n > K_1 \) et  \( y_n \in \mO_2 \) quand \( n > K_2 \). Posons \( K = \max \{ K_1, K_2 \}\). Alors, pour \( n > K \), on a \( (x_n, y_n) \in \mO_1 \times \mO_2 \subset V\). 
\end{proof}

%-----------------------------------------------------------------------------------------------------------------
\subsection{Convergence dans les réels}
%-----------------------------------------------------------------------------------------------------------------

Dans le cas de suites réelles, nous avons la caractérisation suivante qui est souvent donnée comme une définition lorsque seule la topologie sur \( \eR\) est considérée.
\begin{proposition}[Limite d'une suite numérique]	\label{PropLimiteSuiteNum}
	La suite $(x_n)$ est convergente si et seulement s'il existe un réel $\ell$ tel que
	\begin{equation}		\label{EqDefLimSuite}
		\forall \epsilon>0,\,\exists N\in\eN\tq\forall n\geq N,\,| x_n-\ell |<\epsilon.
	\end{equation}
	Dans ce cas, le nombre $\ell$ est la limite de la suite $(x_n)$. Nous dirons aussi souvent que la suite \defe{converge}{} vers le nombre $\ell$.
\end{proposition}
\index{convergence!suite numérique}
\index{limite!suite numérique}

\begin{proof}
  Si \( x_n\to \ell\) et si \( \epsilon>0\) il existe \( N_{\epsilon}\) tel que pour tout \( n\geq N\) nous avons \( x_n\in \mathopen ]\ell-\epsilon,\ell+\epsilon\mathclose[\) (parce que c'est un ouvert contenant \( \ell\)). Mais \( \eR\) est totalement ordonné, et par le lemme~\ref{LemInegEtBoules_CorpsTotalementOrdonne}, cette condition est bien \( | x_n-\ell |<\epsilon\).

    Dans l'autre sens, soit \( \mO\) un ouvert contenant \( \ell\). Par définition de la topologie réelle, il existe \( a, b \in \eR \) tels que \( \ell \in \mathopen]a,b\mathclose[ \subset \mO\). Posons \( \epsilon = \min\{\ell-a, b - \ell \}\): c'est bien un nombre positif. La condition \eqref{EqDefLimSuite} nous assure qu'il existe \( N_{\epsilon} \) tel que pour tout \( n\geq N_{\epsilon}\) nous ayons
    \begin{equation}
        x_n\in \mathopen]\ell-\epsilon, \ell+\epsilon\mathclose[\subset\mO,
    \end{equation}
    ce qui assure que la suite \( (x_n)\) converge vers \( \ell\) pour la topologie usuelle de \( \eR\).
\end{proof}

Une façon équivalente d'exprimer le critère \eqref{EqDefLimSuite} est de dire que pour tout $\epsilon$ positif, il existe un rang $N\in\eR$ tel que l'intervalle $\mathopen[ \ell-\epsilon , \ell+\epsilon \mathclose]$ contient tous les termes $x_n$ au-delà de $N$.

Il est à noter que le rang $N$ dont il est question dans la définition de suite convergente dépend de~$\epsilon$.

Le concept fondamental de la proposition précédente est la notion de valeur absolue qui permet de donner la «distance» entre deux réels. Dans un espace vectoriel normé quelconque, cette notion est généralisée par la distance associée à la norme (définition~\ref{DefEVNetDistance}). Nous pouvons donc facilement obtenir le concept de convergence d'une suite dans un espace vectoriel normé : ce sera la définition~\ref{DefCvSuiteEGVN}.

% On a besoin de cette proposition pour démontrer Weierstrass. Quitte à maintenir, il faut réénoncer pour un espace vectoriel normé et prouver.
\begin{proposition}		\label{PropCvRpComposante}
	Une suite $(x_n)$ dans $\eR^m$ est convergente dans $\eR^m$ si et seulement si les suites de chaque composante sont convergentes dans $\eR$. Dans ce cas nous avons
	 \begin{equation}
		 \lim x_n=\Big( \lim(x_n)_1,\lim (x_n)_2,\ldots,\lim (x_n)_m \Big)
	 \end{equation}
	 où $(x_n)_k$ dénote la $k$-ième composante de $(x_n)$.
\end{proposition}

\begin{example}
	La suite $x_n=\big( \frac{1}{ n },1-\frac{1}{ n } \big)$ converge vers $(0,1)$ dans $\eR^2$. En effet, en utilisant la proposition~\ref{PropCvRpComposante}, nous devons calculer séparément les limites
	\begin{equation}
		\begin{aligned}[]
			\lim\frac{1}{ n }&=0\\
			\lim\big( 1-\frac{1}{ n } \big)&=1.
		\end{aligned}
	\end{equation}
\end{example}

\begin{example}
	Étant donné que la suite $(-1)^n$ n'est pas convergente, la suite $x_n=\big( (-1)^n,\frac{1}{ n } \big)$ n'est pas convergente dans $\eR^2$.
\end{example}

%+++++++++++++++++++++++++++++++++++++++++++++++++++++++++++++++++++++++++++++++++++++++++++++++++++++++++++++++++++++++++++
\section{Connexité}
%+++++++++++++++++++++++++++++++++++++++++++++++++++++++++++++++++++++++++++++++++++++++++++++++++++++++++++++++++++++++++++

%-----------------------------------------------------------------------------------------------------------------
\subsection{Définition et propriétés}
%-----------------------------------------------------------------------------------------------------------------

L'idée de la connexité, c'est de s'assurer qu'un ensemble est «d'un seul tenant».

\begin{definition}  \label{DefIRKNooJJlmiD}
     Lorsque $E$ est un espace topologique, nous disons qu'un sous-ensemble $A$ est \defe{non connexe}{connexité!définition} quand on peut trouver des ouverts $O_1$ et $O_2$ disjoints tels que
    \begin{equation}    \label{EqDefnnCon}
        A=(A\cap O_1)\cup (A\cap O_2),
    \end{equation}
    et tels que $A\cap O_1\neq\emptyset$, et $A\cap O_2\neq\emptyset$. Si un sous-ensemble n'est pas non-connexe, alors on dit qu'il est \defe{connexe}{ensemble connexe}.
\end{definition}
Une autre façon d'exprimer la condition \eqref{EqDefnnCon} est de dire que $A$ n'est pas connexe quand il est contenu dans la réunion de deux ouverts disjoints qui intersectent tous les deux $A$.

\begin{proposition} \label{PropHSjJcIr}
    Soit \( X\) un espace topologique. Les conditions suivantes sont équivalentes.
    \begin{enumerate}
        \item
            L'espace \( X\) est connexe.
        \item
            Si \( X=A\sqcup B\) avec \( A\) et \( B\) fermés disjoints dans \( X\), alors \( A=\emptyset\) ou \( B=\emptyset\).
        \item       \label{ITEMooNIPZooIDPmEf}
            Si \( A\subset X\) avec \( A\) ouvert et fermé en même temps, alors \( A=\emptyset\) ou \( A=X\).
    \end{enumerate}
\end{proposition}
%TODO : une preuve.

Nous verrons plus tard (proposition~\ref{PropConnexiteViaFonction}) une autre caractérisation de la connexité.

\begin{proposition}
    Si \( A\subset X\) est connexe et si \( A\subset B\subset \bar A\), alors \( B\) est connexe.
\end{proposition}
%TODO : une preuve.

\begin{proposition} \label{PropIWIDzzH}
    Stabilité de la connexité par union.
    \begin{enumerate}
        \item
    Une union quelconque de connexes ayant une intersection non vide est connexe.
\item
    Pour tout \( n \in \eN, n > 0 \), si \( A_1,\ldots, A_n\) sont des connexes de \( X\) avec \( A_i\cap A_{i+1}\neq \emptyset\), alors l'union \( \bigcup_{i=1}^nA_i\) est connexe.
    \end{enumerate}
\end{proposition}

\begin{proof}
    Point par point.
    \begin{enumerate}
        \item
    Soient \( \{ C_i \}_{i\in I}\) un ensemble de connexes et un point \( p\) dans l'intersection : \( p\in\bigcap_{i\in I}C_i\). Supposons que l'union ne soit pas connexe. Alors nous considérons \( A\) et \( B\), deux ouverts disjoints recouvrant tous les \( C_i\) et ayant chacun une intersection non vide avec l'union.

    Supposons pour fixer les idées que \( p\in A\) et prenons \( x\in B\cap\bigcup_{i\in I}C_i\). Il existe un \( j\in I\) tel que \( x\in C_j\). Avec tout cela nous avons
    \begin{enumerate}
        \item
            \( C_j\subset A\cup B\) parce que \(A \cup B\) recouvre tous les \( C_i \),
        \item
            \( C_j\cap A\neq \emptyset\) parce que \( p\) est dans l'intersection,
        \item
            \( C_j\cap B\neq\emptyset\) parce que \( x\) est dans cette intersection.
    \end{enumerate}
    Cela contredit le fait que \( C_j\) soit connexe.

\item

    Pour la seconde partie nous procédons de proche en proche\footnote{Parce qu'on a la flemme de faire une preuve par récurrence!}. D'abord \( A_1\cup A_2\) est connexe par la première partie, ensuite \( (A_1\cup A_2)\cup A_3\) est connexe parce que les connexes \( A_1\cup A_2\) et \( A_3\) ont un point d'intersection par hypothèse, et ainsi de suite.
    \end{enumerate}
\end{proof}

\begin{lemma}[Passage de douane\cite{ooDKEWooFqlDyN,ooWBUCooAdPjMK}]        \label{LEMooLKWEooItGnkP}
    Dans un espace topologique, toute partie connexe qui rencontre à la fois une partie \( A\) et son complémentaire rencontre nécessairement la frontière de \( A\).
\end{lemma}

\begin{proof}
    Nommons \( \gamma\) la partie connexe qui intersecte \( A\) et \( A^c\). Les ouverts \( \Int(A)\) et \( X\setminus \bar A\) ne peuvent pas recouvrir \( \gamma\) parce que ce sont deux ouverts disjoints alors que \( \gamma\) est connexe (voir la définition~\ref{DefIRKNooJJlmiD} de la connexité). Donc \( \gamma\) doit contenir des points qui sont dans \( \bar A\) mais pas dans \( \Int(A)\). C'est à dire des points de \( \partial A\).
\end{proof}

%-----------------------------------------------------------------------------------------------------------------
\subsection{Connexité dans \texorpdfstring{$\eR$}{R}}
%-----------------------------------------------------------------------------------------------------------------

Une des nombreuses propositions qui vont servir à prouver le théorème des \href{http://fr.wikipedia.org/wiki/Théorème_des_valeurs_intermédiaires}{valeurs intermédiaires}~\ref{ThoValInter} est la suivante.
\begin{proposition} \label{PropInterssiConn}
    Une partie de $\eR$ est connexe si et seulement si c'est un intervalle.
\end{proposition}
\index{connexité!et intervalles}

\begin{proof}
    La preuve est en deux parties. D'abord nous démontrons que si un sous-ensemble de $\eR$ est connexe, alors c'est un intervalle; et ensuite nous démontrons que tout intervalle est connexe.

    Prenons $A$, une partie de $\eR$ qui n'est pas un intervalle, et prouvons par contraposée qu'il n'est pas connexe. Il existe donc $a$, $b\in A$ et un $x_0$ entre $a$ et $b$ qui n'est pas dans $A$. Comme le but est de prouver que $A$ n'est pas connexe, il faut couper $A$ en deux ouverts disjoints. L'élément $x_0$ qui n'est pas dans $A$ est le bon candidat pour effectuer cette coupure. Prenons $M$, un majorant de $A$ et $m$, un minorant de $A$, et définissons
    \begin{align*}
        \mO_1&=]m,x_0[\\
        \mO_2&=]x_0,M[.
    \end{align*}
    Si $A$ n'a pas de minorant, nous remplaçons la définition de $\mO_1$ par $]-\infty,x_0[$, et si $A$ n'a pas de majorant, nous remplaçons la définition de $\mO_2$ par $]x_0,\infty[$. Dans tous les cas, ce sont deux ensembles ouverts dont l'union recouvre tout $A$. En effet, $\mO_1\cup \mO_2$ contient tous les nombres entre un minorant de $A$ et un majorant sauf $x_0$, mais on sait que $x_0$ n'est pas dans $A$. Cela prouve que $A$ n'est pas connexe.

    Prouvons à présent que tout intervalle est connexe. Pour cela, nous refaisons le coup de \href{http://fr.wikipedia.org/wiki/Contraposée}{la contraposée}. Nous allons donc prendre une partie $A$ de $\eR$, supposer qu'elle n'est pas connexe et puis prouver qu'elle n'est alors pas un intervalle. Nous avons deux ouverts disjoints $\mO_1$ et $\mO_2$ tels que $A\subset \mO_1\cup \mO_2$. Notons \( A_1 = A \cap \mO_1 \) et  \( A_2 = A \cap \mO_2 \); et prenons $a\in A_1$ et $b\in A_2$. Pour fixer les idées, on suppose que $a<b$. Maintenant, le jeu est de montrer qu'il existe un point $x_0$ entre $a$ et $b$ qui ne soit pas dans $A$ (cela montrerait que $A$ n'est pas un intervalle). Nous allons prouver que c'est le cas du point
    \[
      x_0=\sup\{ x\in\mO_1\tq x<b \}.
    \]
    Étant donné que l'ensemble $\mA=\{ x\in\mO_1\tq x<b \}$ est ouvert (c'est l'intersection entre l'ouvert $\mO_1$ et l'ouvert $\mathopen] -\infty, b \mathclose[$, le point $x_0$ n'est pas dans l'ensemble par le lemme~\ref{LemSupOuvPas}. Nous avons donc
    \begin{itemize}
        \item soit $x_0$ n'est pas dans $\mO_1$,
        \item soit $x_0\leq b$,
        \item soit les deux en même temps.
    \end{itemize}
    Nous allons montrer qu'un tel $x_0$ ne peut pas être dans $A$. D'abord, remarquons que $\sup\mA\leq\sup\mO$ parce que $\mA$ est une intersection de $\mO$ avec quelque chose. Ensuite, il n'est pas possible que $x_0$ soit dans $\mO_2$ parce que tout élément de $\mO_2$ possède un voisinage contenu dans $\mO_2$. Un point de $\mO_2$ est donc toujours strictement plus grand que le supremum de $\mO_1$.

    Maintenant, remarquons que si $x_0\leq b$, alors $x_0=b$, sinon $b$ serait un majorant de $\mA$ plus petit que $x_0$, ce qui n'est pas possible vu que $x_0$ est le supremum de $\mA$ et donc le plus petit majorant. Oui mais si $x_0=b$, c'est que $x_0\in\mO_2$, ce qu'on vient de montrer être impossible. Nous voilà déjà débarrassé des deuxièmes et troisièmes possibilités.

    Si la première possibilité est vraie, alors $x_0$ n'est pas dans $A$ parce qu'on a aussi prouvé que $x_0\notin\mO_2$. Or n'être ni dans $\mO_1$ ni dans $\mO_2$ implique de ne pas être dans $A$. Ce point $x_0=\sup\mA$ est donc hors de $A$. Oui, mais comme $a\in\mA$, on a obligatoirement que $x_0\geq a$. Mais par construction, on a aussi que $x_0\leq b$ (ici, l'inégalité est même stricte, mais ce n'est pas important). Donc
    \[
      a\leq x_0\leq b
    \]
    avec $a$, $b\in A$, et $x_0\notin A$. Cela finit de prouver que $A$ n'est pas un intervalle.
\end{proof}


%+++++++++++++++++++++++++++++++++++++++++++++++++++++++++++++++++++++++++++++++++++++++++++++++++++++++++++++++++++++++++++
\section{Compacité}
%+++++++++++++++++++++++++++++++++++++++++++++++++++++++++++++++++++++++++++++++++++++++++++++++++++++++++++++++++++++++++++

La compacité est le thème~\ref{THEMEooQQBHooLcqoKB}.

%---------------------------------------------------------------------------------------------------------------------------
\subsection{Définition et notions connexes}
%---------------------------------------------------------------------------------------------------------------------------

\begin{definition}  \label{DefJJVsEqs}
  Une partie $A$ d'un espace topologique est \defe{compacte}{compact} s'il vérifie la propriété de Borel-Lebesgue : pour tout recouvrement de $A$ par des ouverts (c'est-à-dire une collection d'ouverts dont la réunion contient $A$) on peut tirer un recouvrement fini.
\end{definition}
\begin{remark}
    Certaines sources (dont \wikipedia{fr}{Compacité_(mathématiques)}{wikipédia}) disent que pour être compact il faut aussi être séparé\footnote{Définition~\ref{DefWEOTrVl}.}. Pour ces sources, un espace qui ne vérifie que la propriété de Borel-Lebesgue est alors dit \defe{quasi-compact}{quasi-compact}\index{compact!quasi}.
\end{remark}

\begin{definition}
    Une partie d'un espace topologique est \defe{relativement compact}{compact!relativement}\index{relativement!compact} si son adhérence est compacte.
\end{definition}

\begin{definition}  \label{DefEIBYooAWoESf}
    Un espace topologique est \defe{localement compact}{compact!localement} si tout élément possède un voisinage compact.
\end{definition}

\begin{definition}[Séquentiellement compact]
    Nous disons qu'un espace topologique est \defe{séquentiellement compact}{compact!séquentiellement} si toute suite admet une sous-suite convergente.
\end{definition}

\begin{probleme}
  Étudier les implications - équivalences entre ces définitions.
\end{probleme}

\begin{definition}      \label{DefFCGBooLpnSAK}
    Un espace topologique est \defe{dénombrable à l'infini}{dénombrable!à l'infini} s'il est réunion dénombrable de compacts.
\end{definition}

\begin{definition}
    Une famille \( \mA\) de parties de \( X\) a la \defe{propriété d'intersection finie non vide}{propriété d'intersection non vide} si tout sous-ensemble fini de \( \mA\) a une intersection non vide.
\end{definition}

\begin{proposition}\label{PropXKUMiCj}
    Soient \( X\) un espace topologique et \( K\subset X\). Les propriétés suivantes sont équivalentes :
    \begin{enumerate}
        \item\label{ItemXYmGHFai}
            \( K\) est compact.
        \item\label{ItemXYmGHFaii}
            Si \( \{ F_i \}\) est une famille de fermés telle que \( \bigcap_{i\in I}F_i \cap K =\emptyset\), alors il existe un sous-ensemble fini \( A\) de \( I\) tel que \( \bigcap_{i\in A}F_i \cap K =\emptyset\).
        \item\label{ItemXYmGHFaiii}
            Si \( \{ F_i \}_{i\in I}\) est une famille de fermés telle que pour tout choix de \( A\) fini dans \( I\), \( \bigcap_{i\in A}F_i \bigcap K \neq\emptyset\), alors l'intersection complète est non vide : \( \bigcap_{i\in I}F_i \bigcap K\neq\emptyset\).
        \item\label{ItemXYmGHFaiv}
            Toute famille de fermés de \( X \), à laquelle \( K \) est joint, et qui a la propriété d'intersection finie non vide, a une intersection non vide.
    \end{enumerate}
\end{proposition}

\begin{proof}
    Les propriétés~\ref{ItemXYmGHFaiii} et~\ref{ItemXYmGHFaii} sont équivalentes par contraposition. De plus le point~\ref{ItemXYmGHFaiv} est une simple\footnote{Enfin, simple\dots{} il faut remarquer que dans la formulation de~\ref{ItemXYmGHFaiv}, les intersections peuvent ne pas faire intervenir \( K \), mais, au final, on s'en moque.} reformulation en français de la propriété~\ref{ItemXYmGHFaiii}.

    Prouvons~\ref{ItemXYmGHFai} \( \Rightarrow\)~\ref{ItemXYmGHFaii}. Soit \( \{ F_i \}_{i\in I}\) une famille de fermés tels que \( K\bigcap_{i\in I}F_i=\emptyset\). Les complémentaires \( \mO_i\) de \( F_i\) dans \( X\) recouvrent \( K\) et donc on peut en extraire un sous-recouvrement fini :
    \begin{equation}
        K\subset\bigcup_{i\in A}\mO_i
    \end{equation}
    pour un certain sous-ensemble fini \( A\) de \( I\). Pour ce même choix \( A\), nous avons alors aussi
    \begin{equation}
        \bigcap_{i\in A}F_i \cap K =\emptyset.
    \end{equation}

    L'implication~\ref{ItemXYmGHFaii} \( \Rightarrow\)~\ref{ItemXYmGHFai} est la même histoire de passage aux complémentaires.
\end{proof}

%---------------------------------------------------------------------------------------------------------------------------
\subsection{Quelques propriétés générales}
%---------------------------------------------------------------------------------------------------------------------------

\begin{lemma}   \label{LemOWVooZKndbI}
    Une partie \( K\) d'un espace topologique est compacte si et seulement si de tout recouvrement par des ouverts d'une base de topologie nous pouvons extraire un sous-recouvrement fini.
\end{lemma}
Remarquons que la partie qui est réellement à prouver est que, si \og ça marche \fg{} pour des ouverts d'une base de topologie, alors \og ça marche\fg{} pour tous types d'ouverts.
\begin{proof}
    Soit \( K\) une partie d'un espace topologique et \( \{ \mO_i \}_{i\in I}\) un recouvrement de \( K\) par des ouverts. Chacun des \( \mO_i\) est une union d'éléments de la base de topologie par la proposition~\ref{PropMMKBjgY}: disons \( \mO_i = \bigcup_{j \in J_i} A_{(i,j)} \). Soit \( J = \{ j = (i, j_i) | i \in I, j_i \in J_i \} \); alors nous obtenons  \( \bigcup_{j\in J}A_j=\bigcup_{i\in I}\mO_i\).

    Par hypothèse nous pouvons extraire un ensemble fini \( J_0\subset J\) tel que \( K\subset\bigcup_{j\in J_0}A_j\). Par construction chacun des \( A_j\) est inclus dans (au moins) un des \( \mO_i\). Le choix d'un élément de \( I\) pour chacun des éléments de \( J_0\) donne une partie finie \( I_0\) de \( I\) telle que \( K\subset\bigcup_{j\in J_0}A_j\subset\bigcup_{i\in I_0}\mO_i\).
\end{proof}


\begin{example}[Un compact non fermé]
    En général, un compact n'est pas toujours fermé. Si nous prenons par exemple un ensemble \( X\) de plus de deux points muni de la topologie grossière \( \{ \emptyset,X \}\). Toutes les parties de cet espace sont compactes, mais les seuls fermés sont \( \{ \emptyset,X \}\). Toutes les autres parties sont alors compactes et non fermées.
\end{example}

\begin{proposition}[\cite{OCBrmKo}]\label{PropUCUknHx}
    Tout compact d'un espace topologique séparé est fermé.
\end{proposition}
\index{compact!implique fermé}

\begin{proof}
    Soient \( X\) un espace séparé et \( K\) compact dans \( X\). Nous considérons \( y \in K^c\) et, par hypothèse de séparation, pour chaque \( x\in K\) nous considérons un voisinage ouvert \( V_x\) de \( x\) et un voisinage ouvert~\footnote{Oui, la notation du voisinage peut surprendre, mais elle est quand même pratique pour ce qu'on veut en faire.} \( W_x\) de \( y\) tels que \( V_x\cap W_x=\emptyset\). Bien entendu les \( V_x\) forment un recouvrement ouvert de \( K\) dont nous pouvons extraire un sous-recouvrement fini : soit \( S\) fini dans \( K\) tel que
    \begin{equation}
        K\subset\bigcup_{x\in S}V_x.
    \end{equation}
    L'ensemble \( W=\bigcap_{x\in S}W_x\) est une intersection finie d'ouverts autour de \( y\) et est donc un ouvert autour de \( y\). 
    
    Montrons que \( W\cap K=\emptyset\). Soit \( a\in K\); par définition de \( S\), il existe \( s\in S\) tel que \( a\in V_s\). Par conséquent, \( a\) n'est pas dans \( W_s\) et donc pas non plus dans \( W\).
    

    L'ouvert \( W\) prouve que \( y\) est dans l'intérieur du complémentaire de \( K\), et comme \( y \) est arbitraire, nous concluons que le complémentaire de \( K\) est ouvert (théorème~\ref{ThoPartieOUvpartouv}), en d'autres termes, que \( K\) est fermé.
\end{proof}

\begin{lemma}[\cite{SNSposN}]   \label{LemnAeACf}
    Une partie fermée d'une compact est elle-même compacte.
\end{lemma}
\index{fermé!dans un compact}

\begin{proof}
%    Nous allons utiliser la caractérisation de la compacité en termes de suites donnée par le théorème de Bolzano-Weierstrass~\ref{ThoBWFTXAZNH}. Soit \( K\) un compact et \( F\) un fermé dans \( K\). Nous considérons une suite \( (x_n)\) dans \( F\); par la compacité de \( K\) nous pouvons considérer une sous-suite \( (y_n)\) qui converge dans \( K\) (proposition~\ref{ThoBWFTXAZNH}). Étant donné que \( (y_n)\) est une suite convergente contenue dans \( F\) et étant donné que \( F\) est fermé, la limite est dans \( F\), ce qui prouve que \( (x_n)\) possède une sous-suite convergente dans $F$ et par conséquent que \( F\) est compact.

    Soient \( F\) fermé dans un compact \( K\) et \( \{ \mO_i \}_{i\in I}\) un recouvrement de \( F\) par des ouverts. Vu que \( F\) est fermé, \( F^c\) est ouvert et \( \{ \mO_i \}_{i\in I}\cup\{ K\setminus F \}\) est un recouvrement de \( K\) par des ouverts. Si nous en extrayons un sous-recouvrement fini, c'est un recouvrement de \( F\), et en supprimant éventuellement l'ouvert \( K\setminus F\), ça reste un sous-recouvrement fini de \( F\) tout en étant extrait de \( \{ \mO_i \}_{i\in I}\).
\end{proof}

\begin{proposition}     \label{PropGBZUooRKaOxy}
    Si \( V\) est une partie de l'espace topologique \( X\) muni de la topologie induite\footnote{Définition \ref{DefVLrgWDB}.} \( \tau_V\) de celle de \( X\), et si \( K\) est un compact de \( (V,\tau_V)\) alors \( K\) est un compact de \( (X,\tau_X)\).
\end{proposition}

\begin{proof}
    Soient \(   (\mO_\alpha)_{\alpha\in A}  \) des ouverts de \( X\) recouvrant \( K\). Alors les ensembles \( V\cap \mO_{\alpha}\) recouvrent également \( K\), mais sont des ouverts de \( V\). Donc il en existe un sous-recouvrement fini. Soient donc \( (V\cap\mO_i)_{i\in I}\) recouvrant \( K\) avec \( I\) un sous-ensemble fini de \( A\). Les ensembles \( (\mO_i)_{i\in I}\) recouvrent encore \( K\) et sont des ouverts de \( X\).
\end{proof}

%----------------------------------------------------------------------------------------------------------------------
\subsection{Compacité pour les réels}
%----------------------------------------------------------------------------------------------------------------------


\begin{lemma}[\cite{JUwQXOF}]\label{LemOACGWxV}
    Si \( a<b\in \eR\) alors le segment \( \mathopen[ a , b \mathclose]\) est compact\footnote{Définition~\ref{DefJJVsEqs}}.
\end{lemma}
\index{compact!intervalle \( \mathopen[ a , b \mathclose]\)}

\begin{proof}
    Soit \( \{ \mO_i \}_{i\in I}\) un recouvrement de \( \mathopen[ a , b \mathclose]\) par des ouverts. Nous posons
    \begin{equation}
        M=\{ x\in\mathopen[ a , b \mathclose]\tq \mathopen[ a , x \mathclose] \text{ admet un sous-recouvrement fini extrait de } \{ \mO_i \}_{i\in I} \}.
    \end{equation}
    Notre but est de prouver que \( b\in M\).
    \begin{subproof}

    \item[\( a\) est dans \( M\)]

        Le point \( a\) est naturellement dans un des \( \mO_i\). L'intervalle \( \mathopen[ a , a \mathclose]\) est donc recouvert par un seul des \( \mO_i\).

    \item[\( M\) est un intervalle]

        Soient \( m\in M\) et \( m'\in\mathopen[ a , m [\). Le sous-recouvrement fini qui recouvre \( \mathopen[ a , m \mathclose]\) recouvre a fortiori \( \mathopen[ a , m' \mathclose]\).

    \item[Les trois possibilités restantes]
        À ce niveau de la preuve, il reste trois possibilités pour \( M\) soit il est de la forme \( \mathopen[ a , c \mathclose]\) ou \( \mathopen[ a , c [\) avec \( c<b\), soit il est de la forme \( \mathopen[ a , b \mathclose]\). Nous allons maintenant éliminer les deux premiers cas.

    \item[Ce que \( M\) n'est pas]

        D'abord \( M\) n'est pas de la forme \( \mathopen[ a , c [\) avec \( c<b\). Par l'absurde, commençons par considérer \( \mO_{i_0}\) un ouvert du recouvrement qui contient \( c\); choisissons  \(m \in \mO_{i_0}\) tel que \( m<c\). Alors \( m \in M \), et, si nous joignons \( \mO_{i_0}\) à un recouvrement fini de \( \mathopen[ a , m \mathclose]\) alors nous avons un recouvrement fini de \( \mathopen[ a , c \mathclose]\). On en déduit \( c\in M\).

        Ensuite \( M\) n'est pas de la forme \( \mathopen[ a , c \mathclose]\) avec \( c<b\). En effet si on a un recouvrement fini de \( \mathopen[ a , c \mathclose]\) par des ouverts, alors un de ces ouverts contient \( c\) et donc contient des éléments de \( \mathopen[ a , b \mathclose]\) plus grands que \( c\).
    \end{subproof}
    Nous déduisons que \( M=\mathopen[ a , b \mathclose]\) et qu'il est possible d'extraire un sous-recouvrement fini recouvrant \( \mathopen[ a , b \mathclose]\).
\end{proof}

\begin{lemma}[\cite{MonCerveau}]\label{LemCKBooXkwkte}
    Si \( K_1\) et \( K_2\) sont des compacts dans \( \eR\) alors \( K_1\times K_2\) est compact dans \( \eR^2\).
\end{lemma}

\begin{proof}
    Soit \( \{ \mO_i \}_{i\in I}\) un recouvrement de \( K_1\times K_2\) par des ouverts; grâce au lemme~\ref{LemOWVooZKndbI} nous pouvons supposer que ce sont des carrés. Pour chaque \( x\in K_1\), l'ensemble \( \{ x \}\times K_2\) est compact et donc recouvert par un nombre fini des \( \mO_i\). Soit \( R_x\) un ensemble fini des \( \mO_i\) recouvrant \( \{ x \}\times K_2\).

    Vu que \( R_x\) est une collection finie de carrés nous pouvons considérer \( m_x\), le minimum des rayons. L'ensemble \( K_1\) est recouvert par les boules \( B(x,m_x)\) et il existe donc une collection finie de \( \{ x_i \}_{i\in A}\) tels que \( B(x_i,m_{x_i})\) recouvre \( K_1\).

    Alors \( \{ R_{x_i} \}_{i\in A}\) recouvre \( K_1\times K_2\) parce que \( R_{x_i}\) recouvre l'ensemble \( B(x_i,m_{x_i})\times \{ K_2 \}\).
\end{proof}

%--------------------------------------------------------------------------------------------------------------------------- 
\subsection{Conséquence: les fermés bornés sont compacts}
%---------------------------------------------------------------------------------------------------------------------------

\begin{theorem}[Théorème de Borel-Lebesgue] \label{ThoXTEooxFmdI}
    Une partie de \( \eR^n \) est compacte si et seulement si elle est fermée et bornée.
\end{theorem}
\index{théorème!Borel-Lebesgue}
\index{compact!fermé et borné}

\begin{proof}
    Sens direct.
    \begin{subproof}
    \item[Compact implique borné]
        En effet si \( K\) est non borné dans \( E\) alors \( K\) contient une suite \( (x_n)\) avec \( \| x_n \|>n\). Les boules \( B_i(x_i,\frac{ 1 }{3})\) sont disjointes. On pose \( \mO_0=\complement\bigcup_i\overline{ B(x_i,\frac{1}{ 5 }) }\), qui est ouvert comme complément d'un fermé. Pour \( i\geq 1\) nous posons \( \mO_i=B(x_i,\frac{1}{ 4 })\). Nous avons
        \begin{equation}
            K\subset\bigcup_{i\in \eN}\mO_i
        \end{equation}
        mais vu que \( x_i\) est uniquement dans \( \mO_i\), nous ne pouvons pas extraire de sous-recouvrement fini.
    \item[Compact implique fermé]
        Cela est la proposition~\ref{PropUCUknHx}.
    \end{subproof}
    Sens réciproque.
    \begin{subproof}
    \item[Un intervalle fermé et borné est compact dans \( \eR\)]
        C'est le lemme~\ref{LemOACGWxV}.
    \item[Un produit de segments est compact]
        Le produit de deux compacts de \( \eR\) est un compact dans \( \eR^2\) par le lemme~\ref{LemCKBooXkwkte}.
    \item[Un fermé et borné est compact]
        Soit \( K\) fermé et borné. Vu que \( K\) est borné, il est contenu dans un produit de segments. L'ensemble \( K\) est donc compact parce que fermé dans un compact, lemme~\ref{LemnAeACf}.
    \end{subproof}
\end{proof}

%+++++++++++++++++++++++++++++++++++++++++++++++++++++++++++++++++++++++++++++++++++++++++++++++++++++++++++++++++++++++++++
\section{Limites et continuité de fonctions}
%+++++++++++++++++++++++++++++++++++++++++++++++++++++++++++++++++++++++++++++++++++++++++++++++++++++++++++++++++++++++++++
%---------------------------------------------------------------------------------------------------------------------------
\subsection{Limites}
%---------------------------------------------------------------------------------------------------------------------------

\begin{definition}[Limite d'une fonction, thème~\ref{THEMEooGVCCooHBrNNd}]\label{DefYNVoWBx}
    Soient \( X\) et \( Y\) des espaces topologiques, \( A\subset X\) et \( a\in\bar A\). Soit encore une fonction \( f\colon A\to Y\). L'élément \( y\in Y\) est une \defe{limite}{limite!d'une fonction} de \( f\) en \( a\) si pour tout voisinage \( V\) de \( y\) (pour la topologie de \( Y\)), il existe un voisinage \( W\) de \( a\) dans \( X\) tel que
    \begin{equation}
        f\big( W\cap A\setminus\{a\} \big)\subset V.
    \end{equation}

    Si un tel élément est unique\footnote{Rappelons que ce n'est pas toujours le cas, mais que ça l'est si l'espace topologique est séparé -- définition~\ref{DefYFmfjjm}.}, alors nous disons que cet élément est la \defe{limite}{limite!d'une fonction} de \( f\) et nous notons
    \begin{equation}
        \lim_{x\to a} f(x)=y.
    \end{equation}
\end{definition}

\begin{remark}
    Nous ne saurions trop insister sur le fait que la valeur de \( f\) en \( a\) n'intervient pas dans la définition de la limite de \( f\) en \( a\). Il n'est même pas pas nécessaire que \( f\) soit définie en \( a\) pour que l'on puisse parler de limite de \( f\) en \( a\). Par exemple nous avons
    \begin{equation}
        \lim_{x\to 1} \frac{ x^2-1 }{ x-1 }=2,
    \end{equation}
    alors que la fonction n'est pas définie en \( x=1\).

    Plus généralement, un peu par principe, toutes les fois que la notion de limite apporte une information, le point où l'on prend la limite est spécial. Sinon on ne calculerait pas la limite, mais on regarderait directement la valeur de la fonction. Cela est typiquement le cas lorsque nous verrons les dérivées. En effet, regardons (en faisans du semblant d'anticiper) la définition  \eqref{DEFooOYFZooFWmcAB}. Dans la formule
    \begin{equation}
        f'(a)=\lim_{x\to a} \frac{ f(x)-f(a) }{ x-a },
    \end{equation}
    la fonction sur laquelle nous prenons la limite n'est \emph{jamais} définie en \( x=a\).

    Cela est intimement lié à ce qu'on raconte dans~\ref{SUBSECooVHKCooYRFgrb}.
\end{remark}

\begin{proposition}[Unicité de la limite pour un espace séparé]\label{PropFObayrf}
    Soient \( X\) un espace topologique, \( A\) une partie de \( X\) et \( Y\) un espace topologique séparé\footnote{Définition~\ref{DefYFmfjjm}.}. Nous considérons une fonction \( f\colon A\to Y\). Si \( a\in\bar A\), alors \( f\) admet au plus une limite en \( a\).
\end{proposition}
\index{limite!unicité}

\begin{proof}
    Soient \( y\) et \( y'\) des limites de \( f\) en \( a\), ainsi que des voisinages \( V\) et \( V'\) de \( y\) et \( y'\). Nous prenons également les voisinages \( W\) et \( W'\) correspondants :
    \begin{subequations}
        \begin{numcases}{}
            f(W\cap A)\subset V\\
            f(W'\cap A)\subset V'.
        \end{numcases}
    \end{subequations}
    Quitte à prendre des sous-ensembles nous pouvons supposer que \( W\) et \( W'\) sont ouverts. Il s'ensuit alors que:
    \begin{itemize}
      \item l'ensemble \( W\cap W'\) est un ouvert contenant \( a\) et intersecte donc \( A\);
      \item l'ensemble \( (W\cap W')\cap A\) est donc non vide;
      \item et donc, \( f(W\cap W'\cap A) \) est aussi non vide.
    \end{itemize}
    Mais
    \begin{equation}
            f(W\cap W'\cap A)\subset f(W\cap A)\subset V,
    \end{equation}
    et
    \begin{equation}
            f(W\cap W'\cap A)\subset f(W'\cap A)\subset V',
    \end{equation}
    d'où \( V \) et \( V'\) ont une intersection. Puisque ces ensembles sont arbitraires, nous avons prouvé que tout voisinage de \( y\) et tout voisinage de \( y'\) ont une intersection non vide; étant donné que \( Y\) est séparé, nous devons avoir \( y=y'\).
\end{proof}

%---------------------------------------------------------------------------------------------------------------------------
\subsection{Continuité}
%---------------------------------------------------------------------------------------------------------------------------

\subsubsection{Définitions et propriétés}
%////////////////////////////

La définition suivante est \emph{la} définition de la continuité dans tous les cas.
\begin{definition}[Fonction continue]\label{DefOLNtrxB}
    Une fonction \( f\colon X\to Y\) est \defe{continue au point}{continue!fonction!en un point} \( a\in X\) si \( f(a)\) est une limite\footnote{Définition~\ref{DefYNVoWBx}.} de \( f\) en \( a\).

    Une fonction \( f\colon X\to Y\) est \defe{continue}{continue!fonction entre espaces topologiques} si pour tout ouvert \( \mO\) de \( Y\), l'ensemble
    \begin{equation}      \label{defFminus1ofaset}
      f^{-1}(\mO) = \{ x \in X \tq f(x) \in \mO\}
    \end{equation}
est ouvert dans \( X\).
\end{definition}
La proposition~\ref{PropQZRNpMn} donnera des détails sur ce qu'il se passe lorsque l'espace est métrique.

\begin{theorem} \label{ThoESCaraB}
    Une fonction \( f\colon X\to Y\) est une fonction continue si et seulement si elle est continue en chacun des points de \( X\).
\end{theorem}

\begin{proof}
    En deux parties.
    \begin{subproof}
    \item[Sens direct]
        Nous supposons que \( f\) est une fonction continue. Soient \( a\in X\) et \( V\) un voisinage de \( f(a)\). Nous considérons \( \mO\), un voisinage ouvert de \( f(a)\) contenu dans \( V\); l'ensemble \( f^{-1}(\mO)\) est alors un ouvert contenant \( a\), et l'image de \( f^{-1}(\mO)\) par \( f\) est bien entendu contenue dans \( V\).

    \item[Sens inverse]

        Soit \( \mO\) un ouvert de \( Y\). Pour prouver que \( f^{-1}(\mO)\) est un ouvert de \( X\), nous allons considérer un élément \( a\in f^{-1}(\mO)\) et montrer qu'il existe un voisinage ouvert de \( a\) contenu dans \( f^{-1}(\mO)\); le théorème~\ref{ThoPartieOUvpartouv} nous assurera alors que \( f^{-1}(\mO)\) est ouvert.

        L'ensemble \( \mO\) est un voisinage ouvert de \( f(a)\) parce que \( a\) a été choisi dans \( f^{-1}(\mO)\). Donc la continuité de \( f\) en \( a\) nous assure qu'il existe un voisinage \( W\) de \( a\) tel que \( f(W)\subset\mO\). En prenant un ouvert contenant \( a\) à l'intérieur de \( W\) nous avons un voisinage ouvert de \( a\) contenu dans \( f^{-1}(\mO)\).
    \end{subproof}
\end{proof}

\begin{remark}
    À cause de l'éventuelle non unicité de la limite, deux fonctions continues et égales sur un sous-ensemble dense ne sont pas spécialement égales. Ce sera vrai sur les espaces métriques et plus généralement pour les espaces séparés. Voir l'exemple~\ref{EXooSHKAooZQEVLB} et la proposition~\ref{PropFObayrf}.
\end{remark}

\subsubsection{Continuité séquentielle}
%///////////////////////

\begin{definition}  \label{DefENioICV}
    Si \( X\) et \( Y \) sont deux espaces topologiques, une fonction \( f\colon X\to \eR\) est \defe{séquentiellement continue}{continuité!séquentielle} en un point \( a\) si pour toute suite convergente \( x_n\to a\) dans \( X\) nous avons \( f(x_n)\to f(x)\) dans \( Y\).
\end{definition}

\begin{proposition}[Caractérisation séquentielle de la limite\cite{MonCerveau}] \label{fContEstSeqCont}
  Soient deux espaces topologiques \( X\) et \( Y\) ainsi qu'une fonction \( f\colon X\to Y\). Soit \( a\in X\) et \( \ell\in Y\). Si
  \begin{equation}
    \lim_{x\to a} f(x)=\ell,
  \end{equation}
  alors, pour toute suite \( (x_k) \) telle que \( x_k \to a \), on a
  \begin{equation}
    \lim f(x_k)=\ell.
  \end{equation}
\end{proposition}

\begin{proof}
  Nous considérons une suite \( (x_k)\) qui converge vers \( a\) dans \( X\). Soient \( V\) un voisinage de \( \ell \) et \( W\) un voisinage de \( a\) tels que \( f(W)\subset V\) (définition~\ref{DefYNVoWBx} de la continuité en un point). Par la convergence \( a_k\to a\),  il existe \( N\) tel que pour tout \( k>N\), \( a_k\in W\), et donc tel que \( f(a_k)\in V\), ce qui donne la continuité séquentielle de \( f\).
\end{proof}

\begin{corollary}[Caractérisation séquentielle de la continuité en un point\cite{MonCerveau}]		\label{PropFnContParSuite}
  Une application continue est séquentiellement continue (quels que soient les espaces de départ et d'arrivée).
\end{corollary}

\begin{proof}
  Il suffit de reprendre mot pour mot la preuve précédente en remplaçant \( \ell \) par \( f(a) \).
\end{proof}

\subsubsection{Continuité et application réciproque}
%////////////////////////

\begin{proposition}	\label{PropoInvCompCont}
Soit $f\colon A\subset\eR^n\to B\subset\eR^m$ une bijection continue. Si $A$ est compact, alors $f^{-1}\colon B\to A$ est continue.
\end{proposition}
\index{réciproque!continuité}

\begin{proposition}		\label{PropIntContMOnIvCont}
Soient $I$ un intervalle dans $\eR$ et $f\colon I\to \eR$ une fonction continue strictement monotone. Alors la fonction réciproque $f^{-1}\colon f(I)\to \eR$ est continue sur l'intervalle $f(I)$.
\end{proposition}
\index{réciproque!continuité}

\subsubsection{Homéomorphisme}
%/////////////////

\begin{definition}
    Un \defe{homéomorphisme}{homéomorphisme} est une application bijective continue entre deux espaces topologiques dont la réciproque est continue. Deux espaces topologiques $X$ et $Y$ pour lesquels il existe un homéomorphisme entre $X$ et $Y$, sont dits \defe{isomorphes}{isomorphisme!d'espaces topologiques}.
\end{definition}

%---------------------------------------------------------------------------------------------------------------------------
\subsection{Continuité et topologie induite}
%---------------------------------------------------------------------------------------------------------------------------
\begin{proposition}[\cite{MonCerveau}]     \label{PROPooNPLBooPfmmym}
    Soit une fonction \( f\colon X\to Y\), continue sur l'ouvert \( A\) de \( X\) au sens où elle est continue en chaque point de \( A\). Alors la fonction restriction \( \tilde f\colon A\to Y\) est également continue pour la topologie sur \( A\), induite\footnote{Exemple~\ref{DefVLrgWDB}.} de \( X\).
\end{proposition}

\begin{proof}
    Soit \( a\in A\), et montrons que \( \tilde f\) est continue en \( a\), c'est à dire que \( \tilde f(a)=f(a)\) soit une limite de \( \tilde f\) en \( a\). Soit un voisinage \( V\) de \( \tilde f(a)\) dans \( Y\). Par la continuité de \( f\), nous avons un ouvert \( W\) de \( X\) tel que 
    \begin{equation}
        f\big( W\setminus\{ a \} \big)\subset V.
    \end{equation}
    La partie \( W\cap A\) est un voisinage de \( a\) pour la topologie de \( A\), et vérifie
    \begin{equation}
        f\big( W\cap A\setminus\{ a \} \big)\subset V.
    \end{equation}
    donc \( f(a)\) est une limite de \( \tilde f\) pour \( x\to a\). La fonction \( \tilde f\colon A\to Y\) est continue en chaque point de \( A\).
\end{proof}

Au niveau de la notion de continuité, il n'y a pas trop de changements en passant de \( \eR\) à \( \eQ\) muni de la topologie induite.

\begin{example}     \label{EXooHWIIooYYbfGE}
    Que signifie d'être continue pour une fonction \( f\colon \eQ\to \eR\) ? D'après le théorème~\ref{ThoESCaraB}, il s'agit d'être continue en chaque point de \( \eQ\). Il s'agit donc, par la définition~\ref{DefOLNtrxB} que pour tout \( q\in \eQ\), le nombre \( f(q)\) soit une limite de \( f\) pour \( x\to q\).

    L'espace d'arrivée étant \( \eR\), un voisinage de \( f(q)\) est pris comme une boule de taille \( \epsilon\). La continuité de \( f\) exige qu'il y ait un voisinage \( W\) de \( q\) dans \( \eQ\) tel que pour tout \( q'\in W\) (différent que \( q\)), \( | f(q)-f(q') |<\epsilon\).

    Qu'est-ce qu'un ouvert dans \( \eQ\) ? D'après la définition~\ref{DefVLrgWDB} de la topologie induite, ce sont les ensembles \( \eQ\cap\mO\) avec \( \mO\) ouvert dans \( \eR\). Tout cela pour dire que pour tout \( \epsilon>0\), il doit exister \( \delta>0\) tel que pour tout \( q'\in \eQ\) tel que \( 0<| q-q' |<\delta\), nous ayons \( | f(q)-f(q') |\).

    Bref, c'est exactement le mécanisme usuel de la continuité sur \( \eR\), sauf qu'il faut seulement considérer les rationnels.
\end{example}

\begin{lemma}[Application partielle\cite{MonCerveau}]       \label{LEMooHAODooYSPmvH}
    Soient trois espaces topologiques \( X_1\), \( X_2\) et \( Y\). Nous considérons une fonction continue \( f\colon X_1\times X_2\to Y\) ainsi que \( x_1\in X_1\). Alors l'application
    \begin{equation}
        \begin{aligned}
            g\colon X_2&\to Y \\
            x_2&\mapsto f(x_1,x_2)
        \end{aligned}
    \end{equation}
    est continue.
\end{lemma}

\begin{proof}
    Soit un ouvert \( \mO\) de \( Y\); par hypothèse sur \( f\), la partie \( f^{-1}(\mO)\) est ouverte dans \( X_1\times X_2\). Notre but est de prouver que \( g^{-1}(\mO)\) est un ouvert de \( X_2\). Nous avons :
    \begin{equation}
        g^{-1}(\mO)=\{ x_2\in X_2\tq (x_1,x_2)\in f^{-1}(\mO) \}.
    \end{equation}
    Nous considérons \( x_2\in g^{-1}(\mO)\) et nous prouvons qu'il existe dans \( X_2\) un voisinage de \( x_2\) entièrement contenu dans \( g^{-1}(\mO)\).

    Étant donné que \( (x_1,x_2)\) est dans \( f^{-1}(\mO)\) qui est ouvert, la définition~\ref{DefIINHooAAjTdY} de la topologie sur \( X_1\times X_2\) nous donne des ouverts \( A_1\) dans \( X_1\) et \( A_2\) dans \( X_2\) tels que
    \begin{equation}
        (x_1,x_2)\in A_1\times A_2\subset f^{-1}(\mO).
    \end{equation}

    Nous montrons à présent que \( A_2\subset g^{-1}(\mO)\). Soit \( y_2\in A_2\). Par construction \( (x_1,y_2)\in A_1\times A_2\subset f^{-1}(\mO)\), donc
    \begin{equation}
        g(y_2)=f(x_1,x_2)\in \mO.
    \end{equation}
    Cela termine la démonstration.
\end{proof}

%---------------------------------------------------------------------------------------------------------------------------
\subsection{Continuité et connexité}
%---------------------------------------------------------------------------------------------------------------------------


\begin{proposition} \label{PropConnexiteViaFonction}
  Un espace topologique \( X \) est connexe si et seulement si toute application continue \( X\to \eZ\) est constante.
\end{proposition}

\begin{proposition}\label{PropGWMVzqb}
    L'image d'un ensemble connexe par une fonction continue est connexe.
\end{proposition}

\begin{proof}
    Soit \( f\colon X\to Y\) une application continue entre deux espaces topologiques, et \( E\) une partie connexe de \( X\). Nous devons montrer que \( f(E)\) est connexe dans \( Y\).

    Par l'absurde nous considérons \( A\) et \( B\), deux ouverts de \( Y\) disjoints recouvrant \( f(E)\). Étant donné que \( f\) est continue, les ensembles \( f^{-1}(A)\) et \( f^{-1}(B)\) sont ouverts dans \( X\). De plus ces deux ensembles recouvrent \( E\).

    Si \( x\) est un élément de \( f^{-1}(A)\cap f^{-1}(B)\), alors \( f(x)\in A\cap B\), ce qui est impossible parce que nous avons supposé que \( A\) et \( B\) étaient disjoints. Par conséquent \( f^{-1}(A)\) et \( f^{-1}(B)\) sont deux ouverts disjoints recouvrant \( E\). Contradiction avec la connexité de \( E\). Nous concluons que \( f(E)\) est connexe.
\end{proof}
Une application de ce théorème sera le théorème de valeurs intermédiaires~\ref{ThoValInter}.

\begin{example}
    Les espaces topologiques \( \eR\) et \( \eR^2\) ne sont pas homéomorphes.
\end{example}

\begin{proof}
    Supposons par l'absurde que \( f\colon \eR\to \eR^2\) soit un  homéomorphisme. Nous posons \( E=f\big( \eR\setminus\{ 0 \} \big)\) et \( z_0=f(0)\). Vu que \( f\) est bijective nous avons
    \begin{equation}
        E=\eR^2\setminus\{ z_0 \},
    \end{equation}
    qui est connexe.

    Vu que \( E\) est connexe et que \( f^{-1}\) est continue, la proposition~\ref{PropGWMVzqb} nous dit que \( f^{-1}(E)\) est connexe. Mais par définition, \( f^{-1}(E)=\eR\setminus\{ 0 \}\) qui n'est pas connexe.
\end{proof}

%---------------------------------------------------------------------------------------------------------------------------
\subsection{Continuité et compacité}
%---------------------------------------------------------------------------------------------------------------------------

\begin{theorem}     \label{ThoImCompCotComp}
L'image d'un compact\footnote{Définition~\ref{DefJJVsEqs}.} par une fonction continue est un compact.
\end{theorem}
Dans le cadre des espaces vectoriels normés, ce théorème est démontré en la proposition~\ref{PropContinueCompactBorne}.

\begin{proof}
    Soit $K\subset X$, un ensemble compact, et regardons $f(K)$; en particulier, nous considérons $\Omega$, un recouvrement de $f(K)$ par des ouverts. Nous avons que
    \begin{equation}
        f(K)\subseteq\bigcup_{\mO\in\Omega}\mO.
    \end{equation}
    Par construction, nous avons aussi
    \begin{equation}
        K\subseteq\bigcup_{\mO\in\Omega}f^{-1}(\mO),
    \end{equation}
    en effet, si $x\in K$, alors $f(x)$ est dans un des ouverts de $\Omega$, disons $f(x)\in \mO$, et évidemment, $x\in f^{-1}(\mO)$.  Les $f^{-1}(\mO)$ recouvrent le compact $K$, et donc on peut en choisir un sous-recouvrement fini, c'est à dire un choix de $\{ f^{-1}(\mO_1),\ldots,f^{-1}(\mO_n) \}$ tels que
    \begin{equation}
        K\subseteq \bigcup_{i=1}^nf^{-1}(\mO_i).
    \end{equation}
    Dans ce cas, nous avons que
    \begin{equation}
        f(K)\subseteq\bigcup_{i=1}^n\mO_i,
    \end{equation}
    ce qui prouve la compacité de $f(K)$.
\end{proof}


